% %%
%DIF LATEXDIFF DIFFERENCE FILE
%DIF DEL cwaight_systems_paper_with_revisions2.tex         Tue Nov  3 23:47:57 2026
%DIF ADD cwaight_systems_paper_with_revisions2_secIV.tex   Wed Nov  4 13:22:32 2026
\documentclass[lettersize,journal]{IEEEtran}
%\documentclass[lettersize,journal,draft]{IEEEtran}

\usepackage{amsmath,amsfonts}
\usepackage{algorithmic}
\usepackage{algorithm}
\usepackage{array}
\usepackage[caption=false,font=normalsize,labelfont=sf,textfont=sf]{subfig}
\usepackage{textcomp}
\usepackage{stfloats}
\usepackage{url}
\usepackage{verbatim}
\usepackage{graphicx}
\usepackage{cite}
\hyphenation{op-tical net-works semi-conduc-tor IEEE-Xplore}
\def\BibTeX{{\rm B\kern-.05em{\sc i\kern-.025em b}\kern-.08em
    T\kern-.1667em\lower.7ex\hbox{E}\kern-.125emX}}
\usepackage{balance}
\usepackage{xcolor}
\newtheorem{proposition}{Proposition}
% Revision-round number markers (2026-11-02). Every marked number is listed in
% Queued_Edits/revision_numbers_ledger.md.
\newcommand{\syn}[1]{\textcolor{red}{#1}}            % placeholder data, replaced by the campaign
\newcommand{\rev}[1]{#1}                             % simulation or analysis result, accepted 2026-11-03 (black); rerun if campaign settings change
\newcommand{\open}[1]{\textcolor{magenta}{[#1]}}     % protocol setting not yet decided
% For submission: redefine all three as {#1} (and \open as {}) after the placeholders are gone.
%DIF PREAMBLE EXTENSION ADDED BY LATEXDIFF
%DIF Markup preamble for make_diff.sh (passed to latexdiff --preamble). %DIF PREAMBLE
%DIF Added text: dark green, underlined. (changebar hangs pdflatex on this paper, so no margin bars.) %DIF PREAMBLE
%DIF Deleted text: grey strikeout. %DIF PREAMBLE
%DIF Added equations: dark green. Deleted equations: grey. %DIF PREAMBLE
%DIF New graphics: green frame. Deleted graphics: half size, grey cross. %DIF PREAMBLE
%DIF Red, blue and magenta are left to the paper's \syn, \rev and \open markers. %DIF PREAMBLE
\RequirePackage[normalem]{ulem} %DIF PREAMBLE
\RequirePackage{xcolor} %DIF PREAMBLE
\definecolor{DIFaddbg}{RGB}{204,238,194} %DIF PREAMBLE
\definecolor{DIFaddtx}{RGB}{0,120,45} %DIF PREAMBLE
\definecolor{DIFdeltx}{gray}{0.55} %DIF PREAMBLE
\DeclareRobustCommand{\DIFhl}{\bgroup\markoverwith{\textcolor{DIFaddbg}{\rule[-0.6ex]{2pt}{2.55ex}}}\ULon} %DIF PREAMBLE
\providecommand{\DIFadd}[1]{{\protect\color{DIFaddtx}\uline{#1}}} %DIF PREAMBLE
\providecommand{\DIFdel}[1]{{\protect\color{DIFdeltx}\sout{#1}}} %DIF PREAMBLE
\providecommand{\DIFaddbegin}{\protect\color{DIFaddtx}} %DIF PREAMBLE
\providecommand{\DIFaddend}{\protect\color{black}} %DIF PREAMBLE
\providecommand{\DIFdelbegin}{\protect\color{DIFdeltx}} %DIF PREAMBLE
\providecommand{\DIFdelend}{\protect\color{black}} %DIF PREAMBLE
\providecommand{\DIFmodbegin}{} %DIF PREAMBLE
\providecommand{\DIFmodend}{} %DIF PREAMBLE
\providecommand{\DIFaddFL}[1]{\DIFadd{#1}} %DIF PREAMBLE
\providecommand{\DIFdelFL}[1]{\DIFdel{#1}} %DIF PREAMBLE
\providecommand{\DIFaddbeginFL}{} %DIF PREAMBLE
\providecommand{\DIFaddendFL}{} %DIF PREAMBLE
\providecommand{\DIFdelbeginFL}{} %DIF PREAMBLE
\providecommand{\DIFdelendFL}{} %DIF PREAMBLE
\providecommand{\DIFscaledelfig}{0.5} %DIF PREAMBLE
\RequirePackage{settobox} %DIF PREAMBLE
\RequirePackage{letltxmacro} %DIF PREAMBLE
\newsavebox{\DIFdelgraphicsbox} %DIF PREAMBLE
\newlength{\DIFdelgraphicswidth} %DIF PREAMBLE
\newlength{\DIFdelgraphicsheight} %DIF PREAMBLE
\LetLtxMacro{\DIFOincludegraphics}{\includegraphics} %DIF PREAMBLE
\providecommand{\DIFaddincludegraphics}[2][]{{\color{DIFaddtx}\setlength{\fboxrule}{1.5pt}\fbox{\DIFOincludegraphics[#1]{#2}}}} %DIF PREAMBLE
\providecommand{\DIFdelincludegraphics}[2][]{% %DIF PREAMBLE
\sbox{\DIFdelgraphicsbox}{\DIFOincludegraphics[#1]{#2}}% %DIF PREAMBLE
\settoboxwidth{\DIFdelgraphicswidth}{\DIFdelgraphicsbox} %DIF PREAMBLE
\settoboxtotalheight{\DIFdelgraphicsheight}{\DIFdelgraphicsbox} %DIF PREAMBLE
\scalebox{\DIFscaledelfig}{% %DIF PREAMBLE
\parbox[b]{\DIFdelgraphicswidth}{\usebox{\DIFdelgraphicsbox}\\[-\baselineskip] \rule{\DIFdelgraphicswidth}{0em}}\llap{\resizebox{\DIFdelgraphicswidth}{\DIFdelgraphicsheight}{% %DIF PREAMBLE
\setlength{\unitlength}{\DIFdelgraphicswidth}% %DIF PREAMBLE
\begin{picture}(1,1)% %DIF PREAMBLE
\thicklines\linethickness{2pt} %DIF PREAMBLE
{\color{DIFdeltx}\put(0,0){\framebox(1,1){}}}% %DIF PREAMBLE
{\color{DIFdeltx}\put(0,0){\line( 1,1){1}}}% %DIF PREAMBLE
{\color{DIFdeltx}\put(0,1){\line(1,-1){1}}}% %DIF PREAMBLE
\end{picture}% %DIF PREAMBLE
}\hspace*{3pt}}} %DIF PREAMBLE
} %DIF PREAMBLE
\LetLtxMacro{\DIFOaddbegin}{\DIFaddbegin} %DIF PREAMBLE
\LetLtxMacro{\DIFOaddend}{\DIFaddend} %DIF PREAMBLE
\LetLtxMacro{\DIFOdelbegin}{\DIFdelbegin} %DIF PREAMBLE
\LetLtxMacro{\DIFOdelend}{\DIFdelend} %DIF PREAMBLE
\DeclareRobustCommand{\DIFaddbegin}{\DIFOaddbegin \let\includegraphics\DIFaddincludegraphics} %DIF PREAMBLE
\DeclareRobustCommand{\DIFaddend}{\DIFOaddend \let\includegraphics\DIFOincludegraphics} %DIF PREAMBLE
\DeclareRobustCommand{\DIFdelbegin}{\DIFOdelbegin \let\includegraphics\DIFdelincludegraphics} %DIF PREAMBLE
\DeclareRobustCommand{\DIFdelend}{\DIFOdelend \let\includegraphics\DIFOincludegraphics} %DIF PREAMBLE
\LetLtxMacro{\DIFOaddbeginFL}{\DIFaddbeginFL} %DIF PREAMBLE
\LetLtxMacro{\DIFOaddendFL}{\DIFaddendFL} %DIF PREAMBLE
\LetLtxMacro{\DIFOdelbeginFL}{\DIFdelbeginFL} %DIF PREAMBLE
\LetLtxMacro{\DIFOdelendFL}{\DIFdelendFL} %DIF PREAMBLE
\DeclareRobustCommand{\DIFaddbeginFL}{\DIFOaddbeginFL \let\includegraphics\DIFaddincludegraphics} %DIF PREAMBLE
\DeclareRobustCommand{\DIFaddendFL}{\DIFOaddendFL \let\includegraphics\DIFOincludegraphics} %DIF PREAMBLE
\DeclareRobustCommand{\DIFdelbeginFL}{\DIFOdelbeginFL \let\includegraphics\DIFdelincludegraphics} %DIF PREAMBLE
\DeclareRobustCommand{\DIFdelendFL}{\DIFOdelendFL \let\includegraphics\DIFOincludegraphics} %DIF PREAMBLE
%DIF latexdiff appends its own math-safe \sout after this file, so it is not repeated here. %DIF PREAMBLE
%DIF AMSMATHULEM PREAMBLE %DIF PREAMBLE
\makeatletter %DIF PREAMBLE
\let\sout@orig\sout %DIF PREAMBLE
\renewcommand{\sout}[1]{\ifmmode\text{\sout@orig{\ensuremath{#1}}}\else\sout@orig{#1}\fi} %DIF PREAMBLE
\makeatother %DIF PREAMBLE
%DIF LISTINGS PREAMBLE %DIF PREAMBLE
\RequirePackage{listings} %DIF PREAMBLE
\lstdefinelanguage{DIFcode}{ %DIF PREAMBLE
%DIF DIFCODE_BOLD %DIF PREAMBLE
  % unfortunately \bfseries cannot be combined with ttfamily without extra packages %DIF PREAMBLE
  % also morecomment=[il] is broken as of v1.5b of listings at least %DIF PREAMBLE
  % workaround: plot in white with tiny font %DIF PREAMBLE
  % morecomment=[il]{\%DIF\ <\ }, %DIF PREAMBLE
  moredelim=[il][\color{white}\tiny]{\%DIF\ <\ }, %DIF PREAMBLE
  moredelim=[il][\sffamily\bfseries]{\%DIF\ >\ } %DIF PREAMBLE
} %DIF PREAMBLE
\lstdefinestyle{DIFverbatimstyle}{ %DIF PREAMBLE
	language=DIFcode, %DIF PREAMBLE
	basicstyle=\ttfamily, %DIF PREAMBLE
	columns=fullflexible, %DIF PREAMBLE
	keepspaces=true %DIF PREAMBLE
} %DIF PREAMBLE
\lstnewenvironment{DIFverbatim}[1][]{\lstset{style=DIFverbatimstyle,#1}}{} %DIF PREAMBLE
\lstnewenvironment{DIFverbatim*}[1][]{\lstset{style=DIFverbatimstyle,showspaces=true,#1}}{} %DIF PREAMBLE
\lstset{extendedchars=\true,inputencoding=utf8}

%DIF END PREAMBLE EXTENSION ADDED BY LATEXDIFF

\begin{document}
\title{Adaptive Navigation of Multirobot Systems to Critical Points in 2D Vector Fields in Simulation and Experiment}
\author{Christopher Waight and Christopher A. Kitts,~\IEEEmembership{Senior Member,~IEEE}
\thanks{Manuscript received Month Day, Year; revised Month Day, Year.}
\thanks{C. Waight is a PhD Candidate in the Robotic Systems Laboratory, Department of Mechanical Engineering, Santa Clara University, Santa Clara, CA 95053 USA (e-mail: cwaight@scu.edu)}
\thanks{C. A. Kitts is Professor and Director of the Robotic Systems Laboratory, Department of Mechanical Engineering, Santa Clara University, Santa Clara, CA 95053 USA (e-mail: ckitts@scu.edu)}
}
\markboth{Journal of \LaTeX\ Class Files,~Vol.~XX, No.~X, Month~Year}%
{Waight and Kitts: Adaptive Navigation of Multirobot Systems to Critical Points in 2D Vector Fields}

%\IEEEpubid{0000--0000/00\$00.00~\copyright~2025 IEEE}

\maketitle

\begin{abstract}
    A variety of interesting phenomena exist at the critical points of environmental vector fields. Existing approaches for navigation to critical points either require prior knowledge of the environment or use local information to provide directional guidance only. This paper presents a multirobot method for estimating critical point locations without prior field knowledge, using only real-time local measurements from a minimum of three robots. The sensing is spatially distributed across the robots, while the estimate and control laws are computed on a central computer.
    This estimation enables new control laws to navigate to these points and to orbit them at a fixed radius.
    The estimation error is bounded, and the attraction error stays bounded under actuator lag, sampling, and friction.
    Simulation verifies both attraction and orbital control laws across six canonical linear vector fields and two nonlinear fields using three omnidirectional robots. Hardware validation within a lab-based experimental testbed demonstrates a \syn{100}\% (saddle) and \syn{86}\% (vortex) attraction convergence success rate within 0.05~m across \syn{158} trials, with mean position errors of \syn{0.037}~m or less from the true critical points. An additional \open{36} orbital trials at \open{six} commanded radii held bounded orbits, with radial bias within a quarter of the radius from \syn{0.40}~m (vortex) and \syn{0.50}~m (saddle) upward. % [SYNTHETIC] hardware numbers; [OPEN] orbit trial count and radii
    While these findings establish the estimation method under lab conditions where printed fields do not advect the robots, the present controller does not yet compensate for the disturbance in oceanic and atmospheric flows where the field physically pushes the platform.

\end{abstract}


\begin{IEEEkeywords}
Multirobot systems, adaptive navigation, vector fields, critical points, cooperative estimation, cluster control
\end{IEEEkeywords}


\section{Introduction}

\IEEEPARstart{M}{ultirobot} systems enhance redundancy, throughput, coverage, and reconfigurability in applications ranging from warehouse fulfillment to search and rescue \cite{1}. Equipped with instruments, they form mobile sensor networks that exploit their spatial distribution to estimate and characterize environmental features \cite{2,3,4}. Adaptive navigation (AN) techniques use these estimates to direct the network to or along features of interest. Unlike survey-oriented navigation, which exhaustively traverses a field and identifies features afterward, AN quickly locates local features and follows them in dynamic environments \cite{5,6,7}.

AN in scalar fields is well explored. Each robot collects a scalar measurement as input to a navigation policy that reaches and follows features such as extrema, saddle points, ridges, trenches, contour lines, and fronts \cite{5,6,7,8,9}. Strategies include gradient following \cite{10}, Hessian estimation \cite{4}, extended information consensus filters \cite{11,13}, constrained nonsmooth distributed optimization \cite{14}, and bio-inspired search \cite{15}, and several have been field-deployed on ground, surface, and aerial vehicles for environmental monitoring, pollution tracking, and search and rescue \cite{10,18,19,20}.

By contrast, only a few works address navigation of physically sensed vector fields \cite{21,22,23}. Much of the broader vector field literature constructs artificial fields for path-following and obstacle avoidance, a separate problem from sensing a real flow \cite{24,25,26,28}, and field campaigns that track features carried by an actual ocean current underscore the distinction \cite{29}.

The oceanic and atmospheric environments that motivate this work are governed by a spatially varying velocity or force, and their transport and accumulation behavior is organized by the critical points of the field, the coordinates where the flow vanishes \cite{32,36}. A surface convergence zone acts as a sink, gathering buoyant debris and drifting bodies; this makes such zones high probability search regions after a loss at sea \cite{37,38}. Containing a ruptured wellhead or a harbor discharge requires reaching the source against its own outward flow \cite{40,41}. A mesoscale ocean eddy acts as a center, trapping a body of water and the plankton bloom inside it, and carrying both across a basin for weeks \cite{43,44,45}. Monitoring the bloom over its lifetime therefore requires a sensor to orbit the drifting core at a fixed standoff; a fixed station would not stay with the eddy. These applications need the coordinate of the feature itself, whether to drive the team onto it or to hold it at a fixed standoff radius; a heading toward the feature is not sufficient.

Work that senses the physical field falls into two camps. One camp presumes prior knowledge of the environment, which is unavailable in the unmapped settings that motivate this work \cite{30,31}. The other operates online in two behaviors: minimally actuated drifters that ride and station-keep along ambient streamlines, and PIM-triple-inspired teams that straddle a saddle and track a single stable or unstable manifold from local velocity measurements \cite{23,32,33}. Both supply a heading along the flow or track the manifold of a single assumed saddle; neither recovers the coordinate of a critical point together with its identity across the full taxonomy of feature types. Estimating the locations of critical points in vector fields remains an open problem for multirobot systems.

Recovering those coordinates is harder than following a heading, for reasons of both information and implementation. The information problem is that the location and classification of a critical point are absent from any single velocity sample; they are carried by the spatial derivatives of the field \cite{46}. A single platform must move, sample again, and difference its readings, which conflates the spatial gradient with the field's change in time and can return features that are not physically present. A distributed team removes this temporal aliasing by sampling at one instant \cite{22,23,47}, but recovering the derivatives is still not recovering the feature. Magnitude-based methods \cite{48} reduce the flow to a scalar and yield a descent direction, but the magnitude vanishes at every critical point regardless of type, and estimating the two velocity components as independent scalar fields still does not locate the point. Both the coordinate and the identity are properties of the assembled Jacobian. Inverting it recovers the location where the field vanishes, and its eigenvalues determine the type \cite{36}.

The implementation problem is that the local field must be reconstructed on hardware from noisy distributed measurements. The reconstruction's conditioning degrades as the formation approaches collinearity, so the sampling shape must be held by a formation controller under real actuator dynamics and communication latency. Both the estimator and the cluster's mechatronics limit the accuracy of the recovered features.

This paper introduces a cooperative Jacobian-inversion framework that shifts multirobot vector field navigation from local directional guidance to explicit coordinate inference and feedback control. Fitting a local affine model to three non-collinear measurements is algebraically compact, but closing a cluster control loop on its zero addresses three limitations of existing flow navigation methods. Directional policies \cite{23,33} and magnitude descent \cite{48} guide the robots along a heading or a manifold, which cannot by itself hold a standoff distance. Inverting the assembled Jacobian $\hat{\mathbf{J}}$ instead yields the coordinate $\hat{\mathbf{p}}^*$, which both brings the cluster to rest at the critical point and holds circular orbits at any standoff radius $r_d$. Because the full $2 \times 2$ Jacobian is preserved, one estimator also locates and classifies all six nondegenerate critical point types in a single time step, without the temporal aliasing of a single moving sensor. The paper also links formation size to closed-loop accuracy under real robot dynamics. It proves the three-robot minimum, quantifies the trade-off between field curvature and measurement error, and derives the conditions for sampled-loop stability, the settling error under stiction, and the orbit radius that actuator lag allows. In the present implementation the robots sense and move, and a central computer forms the estimate and computes the commands.

This paper makes four contributions. First, it presents an algorithm that estimates critical point locations from instantaneous, spatially distributed measurements alone. It also shows how field curvature, formation size, and measurement error set the error of the estimate. Second, it proves that three non-collinear robots are the minimum required for this estimation, and that they suffice for every nondegenerate affine field. Third, it demonstrates that this estimate suffices to drive navigation to and around critical points across six vector field types and two nonlinear fields. It also states conditions under which the attraction error stays bounded despite actuator lag, sampled updates, and friction. Fourth, it validates the cooperative estimation and resulting navigation in \open{campaign} hardware experiments that demonstrate repeatable critical point estimation, attraction, and orbit control. % [OPEN] campaign trial count; the submitted 169 is superseded

The remainder of this paper is organized as follows. Section II develops the framework for Jacobian estimation from spatially distributed measurements and critical point prediction, and presents the control laws for approach and orbit. Section III describes the control architecture used to evaluate these strategies in simulation and hardware, and bounds the closed-loop error. Section IV presents simulation results across six vector field types and two nonlinear fields. Section V describes the hardware testbed and its results. Section VI compares the proposed estimate with matched baselines and discusses limitations and future work. Section VII concludes.


\section{Critical Point Detection and Control}

\subsection{Vector Field Environments and Critical Points}

In its most general form, a vector field assigns a vector $\mathbf{v}(\mathbf{p}) \in \mathbb{R}^m$ to each point $\mathbf{p} \in \mathbb{R}^n$ in a given domain, where $m \geq 2$, making vector fields well suited to represent directional physical quantities such as fluid velocity, electromagnetic force, or gravitational force. For the planar fields of this work, this takes the form of a 2D vector $\mathbf{v}(\mathbf{p}) = (u(\mathbf{p}), v(\mathbf{p}))$ assigned to each coordinate $\mathbf{p} = (x, y)$.

Unlike scalar fields where $m=1$, vector fields can exhibit features such as critical points $\mathbf{p}^*$ where the field vanishes, $\mathbf{v}(\mathbf{p}^*) = \mathbf{0}$. Critical points model phenomena such as vortices arising from tidal flows over submarine caves, or plumes when searching for pollution sources \cite{49}.

Near a critical point, the field is well approximated by its linearization, and the local topology is determined by the Jacobian $\mathbf{J} = \nabla\mathbf{v}|_{\mathbf{p}}$. The eigenvalues of $\mathbf{J}$ at $\mathbf{p}^*$ determine the topological type. For $2 \times 2$ systems with nonzero determinant, six standard critical point types arise from the eigenvalue structure, as shown in Fig.~\ref{fig:allfield} and classified in Table~\ref{tab:field_summary}.

\begin{figure}[!t]
    \centering
    \includegraphics[width=\columnwidth]{six_vector_fields_3x2.png}
    \caption{Six linear vector field types with critical points.}
    \label{fig:allfield}
\end{figure}

Eigenvalue analysis can also identify degenerate scenarios such as a zero determinant. In uniform flow fields, for instance, no critical point exists. While these scenarios are detectable through the estimation framework of Section~II-B, they fall outside the scope of the present work. A nonlinear field can also hold several critical points, separated by curves on which the determinant is zero.

\begin{table}[!t]
    \centering
    \caption{Classification of critical points by Jacobian eigenvalue structure ($\lambda$ = eigenvalue; a complex pair $\alpha \pm i\omega$ has real part $\alpha$, imaginary part $\omega$).}
    \label{tab:field_summary}
    \footnotesize
    \setlength{\tabcolsep}{2pt}
    \begin{tabular}{lll}
    \hline
    \textbf{Field Type} & \textbf{Critical Point} & \textbf{Eigenvalues} \\
    \hline
    Vortex & Center & $\pm i\omega$ \\
    Sinking Vortex & Stable Spiral & $\alpha \pm i\omega$, $\alpha < 0$ \\
    Sink & Stable Node & $\lambda_1, \lambda_2 < 0$ \\
    Source & Unstable Node & $\lambda_1, \lambda_2 > 0$ \\
    Spewing Vortex & Unstable Spiral & $\alpha \pm i\omega$, $\alpha > 0$ \\
    Saddle & Saddle Point & $\lambda_1 < 0 < \lambda_2$ \\
    \hline
    \end{tabular}
\end{table}


\subsection{Cooperative Critical Point Estimation}
\noindent

We refer to the group of robots as a cluster: a formation treated as a single virtual rigid body whose centroid, orientation, and internal geometry can be independently specified and controlled \cite{50,51}.

Given a cluster of three robots sampling a planar vector field with robot positions $\mathbf{p}_i = (x_i, y_i)$ and corresponding vector field measurements $(u_i, v_i)$ for $i \in \{1, 2, 3\}$, we can approximate the vector field components across the cluster's spatial extent using linear planar models:
\begin{equation}\label{eq:2}
\hat{U}(x,y) = ax + by + c = u(x,y) + \mathcal{O}(\ell^2)
\end{equation}
\begin{equation}\label{eq:3}
\hat{V}(x,y) = dx + ey + f = v(x,y) + \mathcal{O}(\ell^2)
\end{equation}
where hats denote quantities estimated from the measurements and $\ell$ is the size of the region spanned by the robots and the point $(x,y)$. The $\mathcal{O}(\ell^2)$ term carries the field's curvature, the second-order Taylor terms that a plane cannot represent.

The coefficients are found by solving the linear systems:
\begin{equation}\label{eq:4}
\mathbf{A}\hat{\boldsymbol{\theta}}_u = \mathbf{z}_u, \quad \mathbf{A}\hat{\boldsymbol{\theta}}_v = \mathbf{z}_v
\end{equation}
where:
\begin{equation}\label{eq:5}
\mathbf{A} = \begin{bmatrix} x_1 & y_1 & 1 \\ x_2 & y_2 & 1 \\ x_3 & y_3 & 1 \end{bmatrix}, \quad \hat{\boldsymbol{\theta}}_u = \begin{bmatrix} a \\ b \\ c \end{bmatrix}, \quad \hat{\boldsymbol{\theta}}_v = \begin{bmatrix} d \\ e \\ f \end{bmatrix}
\end{equation}
and $\mathbf{z}_u = [u_1, u_2, u_3]^T$, $\mathbf{z}_v = [v_1, v_2, v_3]^T$ are the collected vector field measurements.

The estimated critical point $\hat{\mathbf{p}}^* = (\hat{x}^*, \hat{y}^*)$ satisfies $\hat{U}(\hat{x}^*, \hat{y}^*) = 0$ and $\hat{V}(\hat{x}^*, \hat{y}^*) = 0$, which in matrix form reads $\hat{\mathbf{J}}\hat{\mathbf{p}}^* = -\hat{\mathbf{h}}$, where the estimated Jacobian matrix and offset vector are:
\begin{equation}\label{eq:8}
\hat{\mathbf{J}} = \begin{bmatrix} a & b \\ d & e \end{bmatrix}, \quad \hat{\mathbf{h}} = \begin{bmatrix} c \\ f \end{bmatrix}
\end{equation}

Thus, provided $\det(\hat{\mathbf{J}}) \neq 0$:
\begin{equation}\label{eq:9}
\hat{\mathbf{p}}^* = -\hat{\mathbf{J}}^{-1}\hat{\mathbf{h}}
\end{equation}

The eigenvalues of $\hat{\mathbf{J}}$ can be used to identify the type of nearby critical point (Table~\ref{tab:field_summary}).

Equation (\ref{eq:9}) returns the zero of the fitted affine field $\hat{\mathbf{v}}(\mathbf{p}) = (\hat{U}(\mathbf{p}), \hat{V}(\mathbf{p})) = \hat{\mathbf{J}}\mathbf{p} + \hat{\mathbf{h}}$, which need not be a zero of the true field. To relate the two, let $\mathbf{p}^*$ be any critical point of the true field, so that $\mathbf{v}(\mathbf{p}^*) = \mathbf{0}$. Substituting $\hat{\mathbf{h}} = -\hat{\mathbf{J}}\hat{\mathbf{p}}^*$ writes the fitted field about its own zero, $\hat{\mathbf{v}}(\mathbf{p}) = \hat{\mathbf{J}}(\mathbf{p} - \hat{\mathbf{p}}^*)$, and evaluating it at $\mathbf{p}^*$ gives
\begin{equation}\label{eq:scope}
\hat{\mathbf{p}}^* - \mathbf{p}^* = -\hat{\mathbf{J}}^{-1}\hat{\mathbf{v}}(\mathbf{p}^*).
\end{equation}
The true field vanishes at $\mathbf{p}^*$, so the term $\hat{\mathbf{v}}(\mathbf{p}^*)$ is the error of the affine model at the true critical point, and the estimate coincides with $\mathbf{p}^*$ exactly when that error is zero. If the true field is affine over a region containing both the cluster and $\mathbf{p}^*$, the models (\ref{eq:2})--(\ref{eq:3}) reproduce it without error and the recovery is exact at any formation size. On a smooth nonlinear field, the fit carries a first-order linearization error that scales with the local curvature of the field; this error decreases as the formation contracts toward the critical point. The inference from $\hat{\mathbf{p}}^*$ to a true critical point is therefore valid only where the affine model remains accurate from the cluster out to $\mathbf{p}^*$.

Both $\hat{\mathbf{J}}$ and $\hat{\mathbf{h}}$ are computed from the same measurements, so their errors are correlated. Because (\ref{eq:scope}) evaluates the fitted field at $\mathbf{p}^*$, the two errors act jointly through the vector $\hat{\mathbf{v}}(\mathbf{p}^*)$ and need not be bounded separately. Taking norms in (\ref{eq:scope}) gives $\|\hat{\mathbf{p}}^* - \mathbf{p}^*\| \le \|\hat{\mathbf{v}}(\mathbf{p}^*)\|/\sigma_{\min}(\hat{\mathbf{J}})$. The singular value $\sigma_{\min}(\hat{\mathbf{J}})$ is the slowest rate at which the fitted field changes in any direction, so a poorly conditioned Jacobian amplifies the model error.

Let $\mathbf{p}_c = \frac{1}{3}\sum_i\mathbf{p}_i$ be the cluster centroid, $\mathbf{e} = \mathbf{p}_c - \mathbf{p}^*$ its displacement from the critical point, and $R$ the distance from $\mathbf{p}_c$ to each robot of the equilateral formation. Let $M$ bound the second derivatives of the field, the largest Hessian norm of $u$ or $v$ over the region holding the robots and $\mathbf{p}^*$, and let $\bar{\varepsilon}$ bound the measurement error of each robot. Expanding the field to second order about $\mathbf{p}_c$ leaves each robot's reading off the linearization by at most $MR^2/\sqrt{2} + \bar{\varepsilon}$. The fit reproduces the linearization exactly and carries these deviations to $\mathbf{p}^*$ through the barycentric weights $w_i = \frac{1}{3} - \frac{2}{3}(\mathbf{p}_i - \mathbf{p}_c)^T\mathbf{e}/R^2$ of $\mathbf{p}^*$, whose magnitudes sum to at most $1 + \sqrt{2}\|\mathbf{e}\|/R$. Adding the linearization's own error $M\|\mathbf{e}\|^2/\sqrt{2}$ at $\mathbf{p}^*$ and bounding the curvature terms by $\tfrac{M}{\sqrt{2}}(R + \|\mathbf{e}\|)^2$ gives
\begin{equation}\label{eq:bound}
\|\hat{\mathbf{p}}^* - \mathbf{p}^*\| \le \frac{\overbrace{\tfrac{M}{\sqrt{2}}(R + \|\mathbf{e}\|)^2}^{\text{curvature}} + \overbrace{\bigl(1 + \tfrac{\sqrt{2}\|\mathbf{e}\|}{R}\bigr)\bar{\varepsilon}}^{\text{measurement error}}}{\sigma_{\min}(\hat{\mathbf{J}})}.
\end{equation}
Formation size moves the two terms in opposite directions. The curvature term grows with $(R + \|\mathbf{e}\|)^2$, the square of the largest distance from a robot to $\mathbf{p}^*$, so a large formation raises it. The measurement term grows with $\|\mathbf{e}\|/R$, the distance the fit extrapolates measured in formation sizes, so a small formation raises it. This extrapolation grows without bound as the robots approach collinearity. The bound measures how far the estimate sits from a true critical point. It assumes one exists nearby and does not establish that it does.

Each robot carries its own sensor, so measurement noise is independent across robots. For zero-mean noise of covariance $\sigma_v^2\mathbf{I}_2$ at each robot on a locally affine field, linearizing (\ref{eq:scope}) in the noise gives
\begin{equation}\label{eq:cov}
\operatorname{Cov}(\hat{\mathbf{p}}^*) \approx \sigma_v^2\, c_w\, \mathbf{J}^{-1}\mathbf{J}^{-T},
\end{equation}
where $c_w = \sum_i w_i^2 = \frac{1}{3} + \frac{2}{3}\|\mathbf{e}\|^2/R^2$ for an equilateral formation. As the cluster reaches the critical point, $c_w \to 1/3$, the variance factor of the mean of three independent readings.


\subsection{Minimum Robot Requirements}

Equations (\ref{eq:2})--(\ref{eq:3}) each contain three coefficients. From (\ref{eq:4}), each robot contributes one equation to each system. Robot $i$ provides $u_i = ax_i + by_i + c$ for the U-component and $v_i = dx_i + ey_i + f$ for the V-component. To uniquely determine three unknowns requires at least three equations. Three robots are therefore the minimum required for critical point estimation.

With fewer than three robots, the system is underdetermined, because the matrix $\mathbf{A}$ in (\ref{eq:4}) has fewer rows than columns. This admits infinitely many solutions for the coefficients, preventing unique identification of the critical point location. With more than three robots ($n > 3$), the system is overdetermined and is solved by least squares.

Three non-collinear robots are also sufficient. Since $\det(\mathbf{A})$ equals twice the signed area of the formation triangle, $\mathbf{A}$ is invertible exactly when the robots are not collinear.

\begin{proposition}[Minimality]\label{prop:min}
For an affine field $\mathbf{v}(\mathbf{p}) = \mathbf{J}(\mathbf{p}-\mathbf{p}^*)$ with $\det(\mathbf{J}) \neq 0$, locating $\mathbf{p}^*$ of every such field from instantaneous measurements requires at least three non-collinear robots, and three suffice.
\end{proposition}

\begin{IEEEproof}
Any configuration of $n \le 2$ robots, or any number of collinear robots, lies on a line $L = \{\mathbf{p}_0 + s\mathbf{d}\}$, with robot $i$ at $s = s_i$. The robots measure $\mathbf{m}_0 + s_i\mathbf{m}_1$, where $\mathbf{m}_0 = \mathbf{J}(\mathbf{p}_0 - \mathbf{p}^*)$ and $\mathbf{m}_1 = \mathbf{J}\mathbf{d}$. Suppose $\mathbf{p}^* \notin L$, and choose any $\mathbf{q} \notin L$. The alternative field $\mathbf{J}_q(\mathbf{p} - \mathbf{q})$ with $\mathbf{J}_q = [\,\mathbf{m}_0\ \ \mathbf{m}_1\,][\,\mathbf{p}_0 - \mathbf{q}\ \ \mathbf{d}\,]^{-1}$ is nonsingular and reproduces every measurement on $L$. Because $\mathbf{q}$ is arbitrary, the measurements cannot determine $\mathbf{p}^*$. Furthermore, $\det(\mathbf{J}_q) = \det(\mathbf{J})\det[\,\mathbf{p}_0 - \mathbf{p}^*\ \ \mathbf{d}\,]/\det[\,\mathbf{p}_0 - \mathbf{q}\ \ \mathbf{d}\,]$ changes sign as $\mathbf{q}$ crosses $L$, so the type of the critical point is also undetermined. Conversely, for three non-collinear robots $\mathbf{A}$ is invertible, so $\hat{\mathbf{J}} = \mathbf{J}$ and $\hat{\mathbf{p}}^* = \mathbf{p}^*$. The same holds for least squares with $n > 3$ robots that are not all collinear.
\end{IEEEproof}


\subsection{Adaptive Navigation Control Laws}
\label{sec:control_primitives}

Two control primitives drive the cluster centroid $\mathbf{p}_c$ toward or around the estimated critical point $\hat{\mathbf{p}}^*$. The mapping from cluster velocity commands to robot velocities is handled by the cluster space controller of Section~III.

\subsubsection{Critical Point Attraction}
To navigate the cluster center $\mathbf{p}_c$ to $\hat{\mathbf{p}}^*$, we apply a proportional control law:
\begin{equation}\label{eq:10}
\dot{\mathbf{p}}_c = k(\hat{\mathbf{p}}^* - \mathbf{p}_c)
\end{equation}
where $k > 0$ is the proportional gain. With $\hat{\mathbf{e}} = \mathbf{p}_c - \hat{\mathbf{p}}^*$ and a stationary estimate, the closed-loop error dynamics are $\dot{\hat{\mathbf{e}}} = -k\hat{\mathbf{e}}$, so $\hat{\mathbf{e}}(t) = \hat{\mathbf{e}}(0)e^{-kt}$. This guarantees global exponential convergence to the estimated critical point with time constant $1/k$, provided the estimate is stationary.

\subsubsection{Orbital Navigation}
To maintain a circular orbit of radius $r_d$ around $\hat{\mathbf{p}}^*$, we define the radial vector $\mathbf{r} = \hat{\mathbf{p}}^* - \mathbf{p}_c$ with magnitude $r = \|\mathbf{r}\|$. The unit radial vector is $\mathbf{n} = \mathbf{r}/r$, and the unit tangent vector $\mathbf{t}$ is obtained by rotating $\mathbf{n}$ by $-90^\circ$. The orbital velocity command is:

\begin{equation}\label{eq:11}
\dot{\mathbf{p}}_c = k_t \mathbf{t} + k_r(r - r_d)\mathbf{n}
\end{equation}

Here, $k_t$ determines the orbital speed (where $\text{sgn}(k_t)$ dictates direction) and $k_r > 0$ drives radial convergence.
This control law decouples radial and angular dynamics. The radial error $e_r = r - r_d$ follows $\dot{e}_r = -k_r e_r$, ensuring exponential convergence to $r_d$. As $r \to r_d$, the angular velocity stabilizes to $\omega = k_t/r_d$, producing steady circular motion.

\section{Multilayer Control Architecture}


The multirobot system control architecture chosen is an established multilayer approach successfully demonstrated in prior cluster space control studies in both 2D \cite{5,8,10} and 3D \cite{7,9} environments. Fig.~\ref{fig:control_architecture} illustrates this three-layered hierarchical structure. The architecture itself is not a contribution of this work. This paper provides the estimation function that fills the feature estimation block and the attraction and orbital controllers that fill the navigation controller block. The vector field measurements enter the architecture only at the adaptive navigation layer, which uses them to form the critical point estimate; the lower layers operate on velocity commands alone. Each layer is discussed briefly.



\subsection{Robot Controller Layer}
\noindent 

Each robot's controller receives velocity commands and translates them into motion. Each robot also samples the local field $\mathbf{v}(\mathbf{p}_i) = (u_i, v_i)$ and reports its global position $\mathbf{p}_i$, providing the inputs to the estimation framework of Section~II.

Details of the hardware platform are described in Section~V. In simulation, each robot is modeled as a single omnidirectional point mass with a first-order response, Coulomb friction, and a speed limit. While the robot moves, its velocity $\mathbf{v}_{r,i} = \dot{\mathbf{p}}_i$ follows its command $\mathbf{v}_{\text{des},i}$ as
\begin{equation}
    \dot{\mathbf{v}}_{r,i}(t) = -\frac{1}{\tau}\mathbf{v}_{r,i}(t) + \frac{1}{\tau}\Bigl(\mathbf{v}_{\text{des},i}(t) - v_c\frac{\mathbf{v}_{r,i}(t)}{\|\mathbf{v}_{r,i}(t)\|}\Bigr),
    \label{eq:continuous}
\end{equation}
where $\tau$ is the actuator time constant \cite{4,7} and the last term is Coulomb friction, an extra command of magnitude $v_c$ opposing the motion. The time constant on the Decabot hardware was found to be approximately 0.28 s.  The robots are also modeled with a maximum speed of $v_{\max} = 0.4$~m/s and a minimum speed of $v_c = 0.025$~m/s needed to overcome stiction (remeasured for this work; an earlier characterization of the testbed reported 0.05~m/s \cite{22}).

For digital implementation at control period $\Delta t$ and cycle $j$, the command is held over each period and friction acts against the frictionless update $\tilde{\mathbf{v}}$, so discretization of (\ref{eq:continuous}) yields
\begin{equation}\label{eq:momentum}
\begin{split}
&\tilde{\mathbf{v}} = \alpha\,\mathbf{v}_{r,i}[j] + (1-\alpha)\,\mathbf{v}_{\text{des},i}[j],\\
&s = \min\bigl\{v_{\max},\, \max\{0,\, \|\tilde{\mathbf{v}}\| - (1-\alpha)v_c\}\bigr\},\\
&\mathbf{v}_{r,i}[j+1] = s\,\tilde{\mathbf{v}}/\|\tilde{\mathbf{v}}\|,
\end{split}
\end{equation}
where $s$ is the new speed and the momentum coefficient $\alpha = e^{-\Delta t/\tau}$. At a 10~Hz control rate, $\alpha \approx 0.70$. Friction removes $(1-\alpha)v_c$ of speed each cycle and stops the robot without reversing it, so a robot at rest moves only once its command exceeds $v_c$. The cap $v_{\max}$ applies last, and the position advances as $\mathbf{p}_i[j+1] = \mathbf{p}_i[j] + \Delta t\,\mathbf{v}_{r,i}[j+1]$. The sensed field enters only the estimate, not the robot dynamics.


\begin{figure}[!t]
    \centering
    \includegraphics[width=\columnwidth]{control_architecture_2.png}
    \caption{Hierarchical control architecture: robot, cluster space, and adaptive navigation layers, as in \cite{5,8}. The contents of the feature estimation and navigation controller blocks are new in this paper.}
    \label{fig:control_architecture}
\end{figure}

\subsection{Cluster Space Controller Layer}

A cluster space controller layer is used to execute centroid velocity commands from the adaptive navigation layer while maintaining formation geometry. The cluster space controller regulates the three-robot triangular formation through a Side-Angle-Side (SAS) parameterization $(p, \beta, q)$, where $p$ is the distance between robots 1 and 2, $\beta$ is the interior angle at robot 2, and $q$ is the distance between robots 2 and 3 \cite{50}. Following the derivation in \cite{50}, differentiating the forward and inverse kinematic mappings between robot coordinates and these cluster variables gives analytic forms for the cluster space Jacobian and its inverse. The inverse Jacobian translates cluster velocity commands into robot velocity commands, and the forward kinematics transform measured robot positions into cluster variables for shape control. The forward kinematics and the cluster space Jacobian close the feedback loop.


\subsection{Adaptive Navigation Layer}

The adaptive navigation layer is the highest level of the control architecture. It specifies cluster mobility and shape based on real time sensor data from each robot using three functional blocks. A feature estimation block computes the critical point estimate $\hat{\mathbf{p}}^*$ using robot positions and vector field measurements (Section~II-B), a navigation controller generates desired centroid velocity commands $\dot{\mathbf{p}}_c$ using the control laws of Section~II-D, and a formation controller regulates the cluster shape variables $(p, \beta, q)$ to desired values using proportional error feedback. The navigation and formation velocity commands are combined into a complete cluster-space velocity vector, which is passed to the cluster space controller layer for transformation into individual robot commands.

\subsection{Closed-Loop Stability and Tracking Bounds}
\label{sec:stability}

Section~II-D assumes a stationary estimate and an ideal centroid, but the estimate changes every cycle and the centroid responds through the sampled lag and Coulomb friction.

\subsubsection{Sampled Loop}
At cycle $j$, write the estimate as $\hat{\mathbf{p}}^*[j] = \mathbf{p}^* + \boldsymbol{\delta}[j]$, so that the attraction law commands $-k(\mathbf{e} - \boldsymbol{\delta})$, where $\mathbf{e}$ is the centroid error of Section~II-B. With the formation held, and apart from friction and the cap, the centroid velocity follows the frictionless update $\tilde{\mathbf{v}}$ of (\ref{eq:momentum}). The $x$ and $y$ axes then decouple, and for each axis $a$ the state $\mathbf{x} = [e_a,\ \dot{p}_{c,a}]^T$ evolves as
\begin{equation}\label{eq:cl}
\mathbf{x}[j{+}1] = \boldsymbol{\Phi}\,\mathbf{x}[j] + \boldsymbol{\Gamma}\,\delta_a[j],
\end{equation}
where
\[
\boldsymbol{\Phi} = \begin{bmatrix} 1{-}(1{-}\alpha)k\Delta t & \alpha\Delta t \\ -(1{-}\alpha)k & \alpha \end{bmatrix}, \quad \boldsymbol{\Gamma} = (1{-}\alpha)k\begin{bmatrix}\Delta t\\ 1\end{bmatrix}.
\]
The Jury test on the characteristic polynomial $z^2 - (1 + \alpha - (1-\alpha)k\Delta t)z + \alpha$ of $\boldsymbol{\Phi}$ gives exponential stability if and only if
\begin{equation}\label{eq:kmax}
0 < k\Delta t < 2(1+\alpha)/(1-\alpha).
\end{equation}
At $\Delta t = 0.1$~s, $\alpha = 0.70$, and $k = 1$~s$^{-1}$, $k\Delta t = 0.1$ against a limit of \rev{11.3}, and both poles have magnitude $\sqrt{\alpha} = \rev{0.84}$. % [REVIEW] 11.3 = 2(1+alpha)/(1-alpha), 0.84 = sqrt(alpha); gamma from stability_note (outside the authorized source list)
A stable linear loop is input-to-state stable, so a bounded estimate error $\boldsymbol{\delta}$, and friction acting as a command disturbance of size at most $v_c$, leave the centroid error bounded.

\subsubsection{Friction and Estimate Error}
Coulomb friction stops the cluster wherever its command is at most $v_c$. The attraction law commands $k\|\mathbf{e} - \boldsymbol{\delta}\| = k\|\mathbf{p}_c - \hat{\mathbf{p}}^*\|$, so the cluster comes to rest inside the friction ring, a disk of radius $v_c/k$ about the estimate, 0.025~m at $k = 1$~s$^{-1}$. The estimate error shrinks as the cluster nears the critical point (\ref{eq:bound}) and settles at the value it has there. Under ideal continuous kinematics, $\dot{\mathbf{e}} = -k(\mathbf{e} - \boldsymbol{\delta})$, and $V = \frac{1}{2}\|\mathbf{e}\|^2$ satisfies
\[
\dot{V} = -k\|\mathbf{e}\|^2 + k\,\mathbf{e}^T\boldsymbol{\delta} \le -k\|\mathbf{e}\|\bigl(\|\mathbf{e}\| - \|\boldsymbol{\delta}\|\bigr),
\]
so the error decreases while the cluster is farther from the critical point than the estimate is. Provided that moving the cluster a given distance changes the estimate error by less than that distance, $\|\partial\boldsymbol{\delta}/\partial\mathbf{p}_c\| < 1$, the changing estimate causes neither instability nor continued hunting, and the cluster comes to rest. It rests within $v_c/k$ of the estimate, and the estimate lies $\|\boldsymbol{\delta}_{\text{final}}\|$ from the critical point, so writing $\mathbf{e}_{\text{final}} = (\mathbf{p}_c - \hat{\mathbf{p}}^*) + \boldsymbol{\delta}_{\text{final}}$ gives $\|\mathbf{e}_{\text{final}}\| \le v_c/k + \|\boldsymbol{\delta}_{\text{final}}\|$.

\subsubsection{Near-Singular Estimates}
The fit matrix $\mathbf{A}$ is singular only if the robots become collinear, which the equilateral formation prevents. The estimate can therefore become singular only through $\hat{\mathbf{J}}$, as in near-uniform flow far from a critical point or on a curve where $\det\mathbf{J} = 0$ between two critical points. When $|\det\hat{\mathbf{J}}|$ falls below a small threshold such as $10^{-10}$, the fit gives no usable estimate, and the controller needs a fallback, such as holding the last valid estimate.
The command limit caps the centroid speed in the meantime. No conditioning gate on $\sigma_{\min}(\hat{\mathbf{J}})$ was applied to (\ref{eq:9}).

\subsubsection{Orbits}
The orbital law (\ref{eq:11}) is undefined at $r = 0$, which the ideal dynamics never reach but a noisy estimate can. When $r < 10^{-6}$~m the implementation sets $\mathbf{n}$ to an arbitrary unit vector, so the command stays defined and its tangential step $k_t\mathbf{t}$ carries the cluster off the singular point on the next cycle.

Lag and friction set the steady orbit radius even with an exact estimate. Friction takes $v_c$ off the orbit speed, $V = k_t - v_c$. In steady circular motion at radius $r^*$ the commanded velocity turns at the rate $V/r^*$, and the lag makes the cluster follow it about $\tau$ late, which tilts its velocity outward by the angle $\tau V/r^*$. The cluster therefore drifts outward at $\tau V^2/r^*$.

To first order, an estimate error adds its radial component to the radial error. A tangential component $\delta_t$ tilts the tangent command by $\delta_t/r_d$, and the radial term cancels the drift $k_t\delta_t/r_d$ at an offset $k_t\delta_t/(k_r r_d)$, largest at small radii.


\section{Control Performance in Simulation}
%DIF >  [REVIEW COPY 2026-11-04] Section IV restructured: the nonlinear fields join IV-A to IV-C, IV-D is cut, and Tables II and III are merged into one table*.
%DIF >  \newnum{} (blue) marks a simulation number that has not been reviewed yet. Rename to \rev{} once approved.
\DIFaddbegin \providecommand{\newnum}[1]{\textcolor{blue}{#1}}
\DIFaddend 

\DIFaddbegin \begin{table*}[!t]
\scriptsize
\centering
\caption{\DIFaddFL{Simulation results on ten fields. Brackets give starts with all three robots inside the $\det\mathbf{J} = 0$ diamond, and the accuracy rows use the trials that succeeded. A dash marks a radius at which most orbits were lost, and $^{a}$ one at which \newnum{48}\% were.}}
\label{tab:simulation_results}\label{tab:simulation_estimation_results}\label{tab:simulation_orbital_results}
\setlength{\tabcolsep}{2.2pt}
\DIFaddFL{\rev{\begin{tabular}{l|c|cc|ccc}
\hline
 & \textbf{Analytical} & \textbf{Vortex} & \textbf{Saddle} & \textbf{Gyre} & \textbf{Curved saddle} & \textbf{Curved saddle} \\
 & \textbf{(6)} & \textbf{(recon.)} & \textbf{(recon.)} & & \textbf{(irrot., $\chi = 1$)} & \textbf{(rot., $\chi = 1$)} \\
\hline
\multicolumn{7}{l}{\textit{(a) Attraction}} \\
\hline
Succ.\ triangle (\%) & 100 & 100 & 100 & \newnum{49 [100]} & \newnum{99.9} & \newnum{72 [100]} \\
Succ.\ any c.p.\ (\%) & 100 & 100 & 100 & \newnum{99.9} & \newnum{99.9} & \newnum{100} \\
Bias (m) & 0.0013 & 0.0205 & 0.0127 & \newnum{0.0018} & \newnum{0.0010} & \newnum{0.0012} \\
$\sigma$ (m) & 0.0249 & 0.0247 & 0.0252 & \newnum{0.0262} & \newnum{0.0242} & \newnum{0.0259} \\
$\bar{e}_f \pm s_e$ (m) & 0.0249 $\pm$ 0.0008 & 0.0293 $\pm$ 0.0132 & 0.0262 $\pm$ 0.0104 & \newnum{0.0259 $\pm$ 0.0044} & \newnum{0.0238 $\pm$ 0.0044} & \newnum{0.0258 $\pm$ 0.0027} \\
Median [p95] (m) & 0.0250 [0.0250] & 0.0321 [0.0500] & 0.0271 [0.0437] & \newnum{0.0249 [0.0334]} & \newnum{0.0231 [0.0309]} & \newnum{0.0253 [0.0299]} \\
Loc.\ error (m) & $<$0.0001 & 0.0222 & 0.0137 & \newnum{0.0057} & \newnum{0.0066} & \newnum{0.0034} \\
$T_{\mathrm{conv}}$ (s) & 2.1 & 2.3 & 2.3 & \newnum{1.8} & \newnum{2.1} & \newnum{1.9} \\
$E_u$ (m$^2$/s) & 0.075 & 0.071 & 0.072 & \newnum{0.047} & \newnum{0.082} & \newnum{0.063} \\
Succ.\ 3/4/5 cm (\%) & 100/100/100 & 45/79/95 & 62/89/100 & \newnum{76/100/100} & \newnum{88/100/100} & \newnum{97/100/100} \\
\hline
\multicolumn{7}{l}{\textit{(b) Orbit: radial bias $\pm$ within-run std (Ctrl.) in m, by $r_d$ (m)}} \\
\hline
\quad 0.01 & +0.129 & +0.133 $\pm$ 0.018 (+0.125) & +0.131 $\pm$ 0.007 (+0.123) & \newnum{+0.127 $\pm$ 0.003 (+0.128)} & \newnum{+0.129 $\pm$ 0.003 (+0.129)} & \newnum{+0.128 $\pm$ 0.002 (+0.129)} \\
\quad 0.10 & +0.093 & +0.094 $\pm$ 0.016 (+0.090) & +0.089 $\pm$ 0.010 (+0.089) & \newnum{+0.080 $\pm$ 0.009 (+0.096)} & \newnum{+0.092 $\pm$ 0.012 (+0.093)} & \newnum{+0.085 $\pm$ 0.006 (+0.095)} \\
\quad 0.20 & +0.067 & +0.089 $\pm$ 0.023 (+0.064) & +0.067 $\pm$ 0.017 (+0.056) & \newnum{+0.023 $\pm$ 0.016 (+0.075)} & \newnum{+0.067 $\pm$ 0.030 (+0.073)} & \newnum{+0.038 $\pm$ 0.011 (+0.073)} \\
\quad 0.30 & +0.051 & +0.087 $\pm$ 0.027 (+0.049) & +0.063 $\pm$ 0.032 (+0.035) & \newnum{$-$0.044 $\pm$ 0.023 (+0.061)} & \newnum{+0.064 $\pm$ 0.063 (+0.075)$^{a}$} & \newnum{$-$0.017 $\pm$ 0.018 (+0.059)} \\
\quad 0.40 & +0.040 & +0.063 $\pm$ 0.030 (+0.040) & +0.040 $\pm$ 0.043 (+0.024) & \newnum{--} & \newnum{--} & \newnum{$-$0.081 $\pm$ 0.024 (+0.050)} \\
\quad 0.50 & +0.033 & +0.026 $\pm$ 0.035 (+0.035) & +0.001 $\pm$ 0.050 (+0.017) & \newnum{--} & \newnum{--} & \newnum{--} \\
\hline
\multicolumn{7}{l}{\textit{(c) Classification (\%)}} \\
\hline
No noise & 100 & 91 & 100 & \newnum{100} & \newnum{100} & \newnum{100} \\
Noise 0.02~m/s & 90--100 & 65 & 100 & \newnum{86} & \newnum{100} & \newnum{100} \\
\hline
\end{tabular}} %DIF >  [REVIEW] experiments/revision_metrics/table_merged.py -> tables/table_sim_merged.tex. Eight-field cells: summary.json and triangle_summary.json (unchanged from the revision's Tables II and III). Nonlinear cells: nonlinear_summary.json (nonlinear.py, Coulomb model, seed 0; noise seed 1)
}\end{table*}

\DIFaddend A Python-based simulation implementing the three-layer control architecture of Section~III was developed to verify the critical point estimation and control primitives in Section~II. \DIFdelbegin \DIFdel{All }\DIFdelend \DIFaddbegin \DIFadd{Unless stated, all }\DIFaddend trials used a cluster of three robots in an equilateral triangular formation with a 0.33~m side length, matching the hardware configuration (Section~V), with control gains of $k = 1$~s$^{-1}$ in (\ref{eq:10}), and $k_t = 0.2$~m/s, $k_r = 0.5$~s$^{-1}$ in (\ref{eq:11}). The robots follow the model of Section~III-A, and the commanded centroid speed was limited to 0.5~m/s.

Six noise-free analytical fields, one for each of the canonical critical point types listed in Table~\ref{tab:field_summary}, were used for trials, each with its critical point at the origin: sink $(-x,-y)$, source $(x,y)$, saddle $(-y,-x)$, vortex $(-y,x)$, sinking vortex $0.4(-y,x)-0.15(x,y)$, and spewing vortex $0.4(y,-x)+0.15(x,y)$. After performance in noise-free environments was verified, simulations were performed on two fields reconstructed from robot-sensed measurements collected on the hardware testbed (vortex and saddle). These reconstructed fields are a fixed spatial distortion of the analytical fields, fit to the robots' readings on the printed maps, and they add no noise that varies in time. Further details on the reconstructed fields are provided in Section~V.

Two nonlinear fields test how the estimate extrapolates. One cell of the double gyre, $\mathbf{v} = 0.1\pi(\cos\pi x \sin\pi y,\allowbreak\ -\sin\pi x \cos\pi y)$, has a center at the origin, saddles at $(\pm 0.5, \pm 0.5)$~m, and $\det\mathbf{J} = 0$ on the diamond $|x| + |y| = 0.5$~m. A curved saddle, $(x, -y)$ at the origin, bends the inflow axis to $v = -y + \chi y^3$, in an irrotational form, $u - iv = z + \chi z^3$ with $z = x + iy$, and a rotational form with stream function $S(x)S(y)$, where $S(s) = \sin(\omega s)/\omega$ and $\omega^2 = 6\chi$. \DIFaddbegin \DIFadd{Both forms use $\chi = 1$ unless stated. The irrotational form then has further saddles at $(0, \pm 1)$~m, and its $\det\mathbf{J}$ vanishes only at the points $(0, \pm 0.58)$~m between them. The rotational form has centers at $(\pm 0.64, \pm 0.64)$~m and $\det\mathbf{J} = 0$ on the diamond $|x| + |y| = 0.64$~m. }\DIFaddend Simulations were therefore run on ten fields in total.

Simulation and hardware trials are scored with the same metrics, and success is geometric.
\begin{itemize}\setlength{\itemsep}{0pt}\setlength{\topsep}{2pt}\setlength{\parsep}{0pt}
\item \emph{Attraction.} A trial ends at time $T$ and succeeds when the critical point lies inside the triangle of the three robots and the centroid never left the map. Its final error is the distance $e_f = \|\mathbf{p}_c(T) - \mathbf{p}^*\|$. The tables give the mean $\bar{e}_f$ and standard deviation $s_e$ of $e_f$ over trials, its median and 95th percentile (p95), and the share of trials with $e_f \le \varepsilon$ for $\varepsilon = 0.03$, $0.04$, and $0.05$~m. Averaging the final positions rather than their distances separates a systematic offset from scatter. The bias is the distance of the mean final position from $\mathbf{p}^*$, and the precision $\sigma$ is the RMS distance of the final positions from that mean, so that bias$^2 + \sigma^2$ is the mean of $e_f^2$. The convergence time $T_{\mathrm{conv}}$ is the median time after which the centroid stays within 0.05~m. Localization error is the mean of $\|\hat{\mathbf{p}}^* - \mathbf{p}^*\|$ over the last 20\% of a trial, and control effort is $E_u = \sum_j \|\mathbf{p}_c[j+1] - \mathbf{p}_c[j]\|^2/\Delta t$. \DIFaddbegin \DIFadd{In a field with several critical points, a trial can succeed at a critical point other than the target $\mathbf{p}^*$. Success there is also scored against the critical point nearest the final centroid, and the metrics above use only the trials that succeeded at $\mathbf{p}^*$.
}\DIFaddend \item \emph{Orbit.} An orbit run is scored on its radial error $\|\mathbf{p}_c - \mathbf{p}^*\| - r_d$ after the first 10~s. Its mean is the radial bias and its standard deviation the within-run spread, and the standard deviation of run means across repeats is the between-trial spread. The controller-only bias (Ctrl.) measures the radius about $\hat{\mathbf{p}}^*$ instead, the radius $r$ that (\ref{eq:11}) regulates, so its gap from the radial bias is the part due to estimation. A run succeeds when the triangle touches the commanded circle of radius $r_d$ about the critical point on at least 95\% of the scored samples and the centroid never left the map. The triangle sets the tolerance, between its inradius of 0.095~m and its circumradius of 0.19~m at the 0.33~m side. A \DIFaddbegin \DIFadd{negative bias means the cluster orbited inside $r_d$. A }\DIFaddend radius is practically achievable when the magnitude of its radial bias is at most $0.25\,r_d$. \DIFaddbegin \DIFadd{A run is lost if the centroid leaves the map or turns less than a quarter revolution about $\mathbf{p}^*$ after the first 10~s. The statistics use the runs that were not lost, and a radius at which most runs were lost is not scored.
}\DIFaddend \item \emph{Classification.} Classification accuracy is the fraction of cycles on which the type of $\hat{\mathbf{J}}$ matches the \DIFdelbegin \DIFdel{true type }\DIFdelend \DIFaddbegin \DIFadd{type of the critical point the trial reached}\DIFaddend .
\end{itemize}

\subsection{Convergence to Critical Points}

To test convergence to critical points, the following simulation experiments were performed. For each of the \DIFdelbegin \DIFdel{eight }\DIFdelend \DIFaddbegin \DIFadd{ten }\DIFaddend fields, the cluster was initialized at a random position drawn uniformly from a 1~m $\times$ 1~m box centered on the \DIFdelbegin \DIFdel{true }\DIFdelend \DIFaddbegin \DIFadd{target }\DIFaddend critical point and the simulation was run for 100 time steps at 10~Hz. This was performed 1000 times per field, \DIFdelbegin \DIFdel{and }\DIFdelend \DIFaddbegin \DIFadd{from the same starts and headings in every field, and }\DIFaddend the final centroid position in each trial was used to determine the metrics. \DIFaddbegin \DIFadd{In the gyre and the rotational saddle, the starts with all three robots inside the $\det\mathbf{J} = 0$ diamond around the target are also scored on their own.
}\DIFaddend 

 The cluster converged to the critical point in \rev{100}\% of 1000 trials for every analytical field at every threshold (Table~\DIFdelbegin \DIFdel{\ref{tab:simulation_estimation_results}}\DIFdelend \DIFaddbegin \DIFadd{\ref{tab:simulation_results}}\DIFaddend ). There the estimate is exact, and every trial parks on the friction ring of Section~III-D (Fig.~\ref{fig:sim_convergence}). The reconstructed fields' fixed distortion moves the rest point off that ring and lowers success at 0.03~m, while at least \rev{95}\% of trials still end within 0.05~m. % [REVIEW] revision_metrics/tables/table_II_revised.tex; triangle rule: revision_metrics/results/triangle_*.csv

With lag and friction off, white noise up to 0.03~m/s changes the reconstructed fields' final error by at most \rev{1.7}~mm, so their error is fixed distortion. On the analytical fields the covariance (\ref{eq:cov}), propagated through the sampled loop, predicts the closed-loop mean error within \rev{4}\% wherever the field changes across the formation by at least three times the noise. % [REVIEW] hyperparameter study v3 E6: median ratio 1.004, range 0.98 to 1.04, over 66 noise settings

\begin{figure}[!t]
    \centering
    \includegraphics[width=\columnwidth]{fig_sim_convergence.pdf}
    \caption{Centroid distance to the critical point in simulated plane-fit attraction, \rev{six} trials per field. On the analytical vortex every trial parks on the friction ring $v_c/k$, and the reconstructed fields' distortion moves the rest point off it.} % [REVIEW] hyperparameter study v3 E1hw saved trajectories
    \label{fig:sim_convergence}
\end{figure}

\DIFdelbegin %DIFDELCMD < \begin{table}[!t]
%DIFDELCMD < \footnotesize
%DIFDELCMD < %%%
\DIFdelendFL \DIFaddbeginFL \DIFaddFL{In the nonlinear fields, the trials that reached the target settled as tightly as in the analytical fields, and the curvature appeared in the localization error. Every successful trial ended within \newnum{0.04}~m of $\mathbf{p}^*$, with a bias of at most \newnum{1.8}~mm. The estimate, however, settled \newnum{5.7}~mm from the gyre's center and \newnum{6.6} and \newnum{3.4}~mm from the irrotational and rotational saddles, against less than \rev{0.1}~mm in the analytical fields. The field's cubic terms do not cancel over an equilateral triangle, so this error grows roughly as the cube of the formation side. In the gyre it was \newnum{0.16}, \newnum{1.2}, \newnum{5.7}, and \newnum{25}~mm at sides of 0.10, 0.20, 0.33, and 0.50~m, and the curved saddles followed the same trend. The formation side changed the success rates by at most \newnum{2.3} points. %DIF >  [REVIEW] nonlinear.py: nonlinear_summary.json "success" blocks and "size" (mean localization error over successful trials; gyre 0.16, 1.15, 5.72, 24.98 mm; irrot 0.09, 1.33, 6.59, 23.7; rot 0.10, 0.72, 3.43, 13.55). Success rates by side: gyre 50.1 to 49.0, irrot 100 to 97.7, rot 73.6 to 72.1
}

\DIFaddFL{From the full box, the cluster reached the target in \newnum{49}\% of gyre trials and \newnum{72}\% of rotational-saddle trials. In all but \newnum{two} of the 3000 nonlinear trials, the cluster ended with some critical point inside its triangle. The attraction law equals $-k\hat{\mathbf{J}}^{-1}\hat{\mathbf{v}}(\mathbf{p}_c)$, a Newton step toward the zero of the local fit, and in continuous time such steps cannot carry the cluster across a curve where $\det\hat{\mathbf{J}} = 0$. The sign of $\det\hat{\mathbf{J}}$ on the first cycle therefore decided the type of critical point reached in all \newnum{2000} gyre and rotational-saddle trials, positive for a center and negative for a saddle (Fig.~\ref{fig:nonlinear}(a)). From starts with all three robots inside the diamond, the cluster reached the target in every trial, \newnum{154} in the gyre and \newnum{334} in the rotational saddle. With the centroid inside but a robot outside, this fell to \newnum{85}\% and \newnum{89}\%, and all but one of the failures began with the wrong sign of $\det\hat{\mathbf{J}}$. Every gyre start within 0.3~m of the center reached it, falling to \newnum{76}\% and \newnum{18}\% in the next two 0.1~m bands, and every rotational-saddle start within 0.4~m reached the target. In the irrotational saddle, whose $\det\mathbf{J}$ is negative everywhere except at two points, every trial ended nearest the target. %DIF >  [REVIEW] nonlinear.py: 1000 starts per field, same draws as sim.py. Whole triangle inside the diamond: gyre 154/154, rotational 334/334 (summary block formation_in_region). Centroid inside, a robot outside: gyre 292/343 (50 wrong-sign, 1 left the map), rotational 361/405 (44 wrong-sign) (block straddle). Centroid inside, any robots: gyre 446/497, rotational 695/739. Start bands from nonlinear_summary.json by_start_distance. Superseded: the 5000-start +-0.45 m gyre run (99.88%, 92.3%)
}

\begin{figure}[!t]
    \DIFaddendFL \centering
    \DIFaddbeginFL \DIFaddFL{\includegraphics[width=\columnwidth]{fig_nonlinear_triangles.pdf}
    }\DIFaddendFL \caption{\DIFdelbeginFL \DIFdelFL{Simulation results across eight field types }\DIFdelendFL \DIFaddbeginFL \DIFaddFL{The gyre cell in simulation, with streamlines in gray, the $\det\mathbf{J} = 0$ diamond dashed, and the formation drawn to scale. }\DIFaddendFL (\DIFdelbeginFL \DIFdelFL{six analytical + two reconstructed fields}\DIFdelendFL \DIFaddbeginFL \DIFaddFL{a}\DIFaddendFL ) \DIFdelbeginFL \DIFdelFL{, }\DIFdelendFL \DIFaddbeginFL \DIFaddFL{The }\DIFaddendFL 1000 \DIFdelbeginFL \DIFdelFL{trials }\DIFdelendFL \DIFaddbeginFL \DIFaddFL{attraction starts }\DIFaddendFL of \DIFdelbeginFL \DIFdelFL{10}\DIFdelendFL \DIFaddbeginFL \DIFaddFL{Table}\DIFaddendFL ~\DIFdelbeginFL \DIFdelFL{s per field. The six analytical fields give identical results except in classification}\DIFdelendFL \DIFaddbeginFL \DIFaddFL{\ref{tab:simulation_results}}\DIFaddendFL , \DIFdelbeginFL \DIFdelFL{and p95 is }\DIFdelendFL \DIFaddbeginFL \DIFaddFL{blue where }\DIFaddendFL the \DIFdelbeginFL \DIFdelFL{95th percentile of $e_f$}\DIFdelendFL \DIFaddbeginFL \DIFaddFL{cluster reached the center and orange where it reached a corner saddle (squares)}\DIFaddendFL . \DIFdelbeginFL \DIFdelFL{Classification uses }\DIFdelendFL \DIFaddbeginFL \DIFaddFL{Both drawn trials start inside }\DIFaddendFL the \DIFdelbeginFL \DIFdelFL{last 2}\DIFdelendFL \DIFaddbeginFL \DIFaddFL{diamond, but the orange one has a robot outside it. (b) Orbits at $r_d = 0.3$}\DIFaddendFL ~\DIFdelbeginFL \DIFdelFL{s of each trial}\DIFdelendFL \DIFaddbeginFL \DIFaddFL{m}\DIFaddendFL , \DIFdelbeginFL \DIFdelFL{without noise }\DIFdelendFL \DIFaddbeginFL \DIFaddFL{which stays inside the diamond }\DIFaddendFL and \DIFdelbeginFL \DIFdelFL{with white measurement noise of 0.02}\DIFdelendFL \DIFaddbeginFL \DIFaddFL{inside $r_d$ (dotted), and at $r_d = 0.4$}\DIFaddendFL ~m\DIFdelbeginFL \DIFdelFL{/s}\DIFdelendFL \DIFaddbeginFL \DIFaddFL{, which reaches the diamond and is lost}\DIFaddendFL .} \DIFdelbeginFL %DIFDELCMD < \label{tab:simulation_estimation_results}
%DIFDELCMD < \setlength{\tabcolsep}{2pt}
%DIFDELCMD < %%%
\DIFdelFL{\rev{\begin{tabular}{lccc}
\hline
\textbf{Metric} & \textbf{Analytical (6)} & \textbf{Vortex (recon.)} & \textbf{Saddle (recon.)} \\
\hline
Bias (m) & 0.0013 & 0.0205 & 0.0127 \\
$\sigma$ (m) & 0.0249 & 0.0247 & 0.0252 \\
$\bar{e}_f \pm s_e$ (m) & 0.0249 $\pm$ 0.0008 & 0.0293 $\pm$ 0.0132 & 0.0262 $\pm$ 0.0104 \\
Median [p95] (m) & 0.0250 [0.0250] & 0.0321 [0.0500] & 0.0271 [0.0437] \\
Loc.\ error (m) & $<$0.0001 & 0.0222 & 0.0137 \\
$T_{\mathrm{conv}}$ (s) & 2.1 & 2.3 & 2.3 \\
$E_u$ (m$^2$/s) & 0.075 & 0.071 & 0.072 \\
Succ.\ 3/4/5 cm (\%) & 100/100/100 & 45/79/95 & 62/89/100 \\
Succ.\ triangle (\%) & 100 & 100 & 100 \\
Class.\ (\%), no noise & 100 & 91 & 100 \\
Class.\ (\%), noise & 90--100 & 65 & 100 \\
\hline
\end{tabular}} %DIF <  [REVIEW] revision_metrics/tables/table_II_revised.tex (six identical analytical rows merged, transposed); classification rows: table_classification_revised.tex, attraction columns
}%DIFDELCMD < \end{table}
%DIFDELCMD < %%%
\DIFdelend %DIF >  [REVIEW] experiments/revision_metrics/fig_nonlinear.py (Coulomb model; trials from nonlinear.py)
    \DIFaddbegin \label{fig:nonlinear}
\end{figure}
\DIFaddend 

\subsection{Orbital Control}

Orbital control was evaluated on the same \DIFdelbegin \DIFdel{eight }\DIFdelend \DIFaddbegin \DIFadd{ten }\DIFaddend fields at different commanded radii. For each radius, the cluster was positioned at the desired radius to initialize the run, then the orbital controller (\ref{eq:11}) was used for 600 time steps. The mean radial error and standard deviation were computed after discarding the first 100 time steps (10~s) as a settling transient. Each run was repeated \rev{100} times per field and radius. % [REVIEW] revision_metrics/tables/table_III_revised.tex caption (100 trials per field and radius)


    Table~\DIFdelbegin \DIFdel{\ref{tab:simulation_orbital_results} }\DIFdelend \DIFaddbegin \DIFadd{\ref{tab:simulation_results}(b) }\DIFaddend presents the orbital tracking results. The six noise-free fields produced identical performance at each radius, with standard deviations below \rev{0.001}~m. This is expected\DIFdelbegin \DIFdel{: on }\DIFdelend \DIFaddbegin \DIFadd{. On }\DIFaddend an exactly linear noise-free field the three-robot fit recovers the critical point exactly at every step, so the controller input, and hence the trajectory, is independent of the field type. \DIFdelbegin \DIFdel{These affine fields therefore do not show how the estimate extrapolates to a nonlinear field. }\DIFdelend Mean radial errors had a positive bias, and decreased with increasing radius. The reconstructed fields repeat the trend, with error standard deviations that increase with the commanded radius. By the $0.25\,r_d$ criterion the practical range is $r_d \ge \rev{0.3}$~m on the analytical fields and the reconstructed saddle, and $r_d \ge \rev{0.4}$~m on the reconstructed vortex. By the triangle rule, no orbit succeeds at $r_d = 0.01$~m and every analytical orbit succeeds from 0.1~m up. The reconstructed saddle succeeds in \rev{43} to \rev{100}\% of runs from 0.1~m up, and the reconstructed vortex in \rev{0}\% at 0.1 and 0.2~m rising to \rev{100}\% at 0.5~m. % [REVIEW] revision_metrics/tables/table_III_revised.tex; triangle rule: revision_metrics/results/triangle_*.csv

\DIFdelbegin %DIFDELCMD < \begin{table}[!t]
%DIFDELCMD <     \centering
%DIFDELCMD <     %%%
%DIFDELCMD < \caption{%
{%DIFAUXCMD
\DIFdelFL{Simulated orbital control across eight field types, radial errors in metres. Bias is $\pm$ its within-run standard deviation, and Ctrl.\ is the controller-only bias. Negative error indicates the cluster orbited inside the commanded radius.}}
    %DIFAUXCMD
%DIFDELCMD < \label{tab:simulation_orbital_results}
%DIFDELCMD <     \footnotesize
%DIFDELCMD <     \setlength{\tabcolsep}{2pt}
%DIFDELCMD <     %%%
\DIFdelFL{\rev{\begin{tabular}{c|c|cc|cc}
    \hline
    \textbf{$r_d$} & \textbf{Noise-free (6)} & \multicolumn{2}{c|}{\textbf{Vortex (recon.)}} & \multicolumn{2}{c}{\textbf{Saddle (recon.)}} \\
    \textbf{(m)} & Bias & Bias $\pm$ std & Ctrl. & Bias $\pm$ std & Ctrl. \\
    \hline
    0.01 & +0.129 & +0.133 $\pm$ 0.018 & +0.125 & +0.131 $\pm$ 0.007 & +0.123 \\
    0.10 & +0.093 & +0.094 $\pm$ 0.016 & +0.090 & +0.089 $\pm$ 0.010 & +0.089 \\
    0.20 & +0.067 & +0.089 $\pm$ 0.023 & +0.064 & +0.067 $\pm$ 0.017 & +0.056 \\
    0.30 & +0.051 & +0.087 $\pm$ 0.027 & +0.049 & +0.063 $\pm$ 0.032 & +0.035 \\
    0.40 & +0.040 & +0.063 $\pm$ 0.030 & +0.040 & +0.040 $\pm$ 0.043 & +0.024 \\
    0.50 & +0.033 & +0.026 $\pm$ 0.035 & +0.035 & +0.001 $\pm$ 0.050 & +0.017 \\
    \hline
    \end{tabular}} %DIF <  [REVIEW] revision_metrics/tables/table_III_revised.tex (RMSE columns omitted; linear std < 0.001 and Ctrl. = Bias)
    }%DIFDELCMD < \end{table}
%DIFDELCMD < %%%
\DIFdelend \DIFaddbegin \DIFadd{In the nonlinear fields, the curvature moved the orbit inward or outward according to how the field's speed grows with distance from $\mathbf{p}^*$. The gyre's speed grows more slowly than in proportion to that distance, so the fit places the center beyond its true position. The cluster holds $r_d$ about this estimate and therefore orbits inside $r_d$, by \newnum{0.044}~m at $r_d = 0.3$~m, while its controller-only bias stays within \newnum{0.01}~m of the analytical bias. The rotational saddle behaves the same way, by up to \newnum{0.081}~m at 0.4~m. In the irrotational saddle the speed grows faster than linearly along the outflow axis and more slowly along the inflow axis. The two shifts cancel in the mean radius, which matches the analytical bias up to 0.2~m, while the within-run spread grows to \newnum{0.030}~m. Across $\chi$ from $-1$ to 1, the rotational form moved the mean radius by up to \newnum{0.17}~m, outward for $\chi < 0$, where its speed grows faster than linearly, and inward for $\chi > 0$. The irrotational form moved it by at most \newnum{0.02}~m wherever most runs were kept. By the $0.25\,r_d$ criterion the practical range is \newnum{0.2} to \newnum{0.3}~m in the gyre and \newnum{0.2} to \newnum{0.4}~m in the rotational saddle, where the inward shift offsets part of the outward drift from lag. %DIF >  [REVIEW] nonlinear.py: nonlinear_summary.json "orbit" and "sweep" (chi in -1, -0.5, 0.5, 1, both forms, six radii, 100 runs each); largest shift rotational chi = -0.5 at r_d 0.4 (+0.214 against +0.040); irrotational chi = -1 at r_d 0.3 (+0.070 against +0.051, 59% kept)
}\DIFaddend 

\DIFaddbegin \DIFadd{At $\chi = 1$, orbits were lost once the formation reached the set on which $\det\mathbf{J} = 0$. Every lost gyre orbit had first crossed the diamond, and all but one ended nearest a corner saddle. These were \newnum{62}\% of runs at $r_d = 0.4$~m and \newnum{93}\% at 0.5~m (Fig.~\ref{fig:nonlinear}(b)). In the rotational saddle \newnum{52}\% of runs at 0.5~m were lost. In the irrotational saddle \newnum{48}\% of runs at 0.3~m and all from 0.4~m were lost. Most left the map along the inflow axis, past the points $(0, \pm 0.58)$~m where $\det\mathbf{J} = 0$ and toward the saddles at $(0, \pm 1)$~m. With $\chi = -1$, the rotational form has no such set, and its outward shift carried every orbit from 0.4~m off the map. %DIF >  [REVIEW] nonlinear.py keep rule (left the map, or < 0.25 revolutions about p* after 10 s). Gyre: 62/62 and 93/93 lost runs crossed |x|+|y| = 0.5 (13 and 7 kept runs also crossed); irrotational chi = 1 at r_d 0.3: 37 of 48 lost runs left the map, median exit 78 deg from the x axis
}

\DIFaddend \subsection{Classification}

A center ($\operatorname{tr}\mathbf{J} = 0$) and the canonical sink and source ($\mathbf{J} = \mp\mathbf{I}$, with repeated eigenvalues) each lie on a boundary between types in Table~\ref{tab:field_summary}, and any error in $\hat{\mathbf{J}}$ can move them across it. The classifier therefore counts an eigenvalue pair within about $6^\circ$ of the imaginary or real axis as a center or a node. Table~\DIFdelbegin \DIFdel{\ref{tab:simulation_estimation_results} }\DIFdelend \DIFaddbegin \DIFadd{\ref{tab:simulation_results}(c) }\DIFaddend gives its accuracy \DIFaddbegin \DIFadd{over the last 2~s of each attraction trial}\DIFaddend . Without noise every field except the reconstructed vortex is labeled correctly on every cycle. With white measurement noise of 0.02~m/s on every reading, the boundary types are labeled least reliably. The center is labeled correctly on \rev{90}\% of cycles and the star nodes on \rev{95}\%, or \rev{100}\% in coarse classes that merge nodes and spirals of the same stability, against \rev{100}\% for the saddle. % [REVIEW] revision_metrics/tables/table_classification_revised.tex, error note

\DIFdelbegin \subsection{\DIFdel{Curvature and Multiple Critical Points}}
%DIFAUXCMD
\addtocounter{subsection}{-1}%DIFAUXCMD
\DIFdelend \DIFaddbegin \DIFadd{In the nonlinear fields, every successful trial labeled the target's type on every cycle without noise, and every trial that ended at another critical point labeled that point's type on every cycle. With noise, the gyre's center was labeled correctly on \newnum{86}\% of cycles, against \rev{90}\% for the analytical center, and both curved saddles on \newnum{100}\%. %DIF >  [REVIEW] nonlinear.py: success blocks cls_target and cls_target_noise; trials at another critical point (gyre 506, rotational 279) any_cls_acc = 1.0, no noise
}\DIFaddend 

\DIFdelbegin \DIFdel{Formation size sets the curvature error (Fig.~\ref{fig:nonlinear}). With lag and friction off, the estimate's fixed point in the gyre lies off the center by an amount growing as the cube of the formation side, because the field's cubic terms do not cancel over an equilateral triangle. At the 0.33~m side it is \rev{6.6}~mm. In the curved saddle, the cubic shifts the estimate along both axes. The shifts cancel in the irrotational form and add in the rotational form, where they move the mean orbit radius by up to \rev{0.18}~m for $|\chi| \le 1$. Attraction still ended within 0.05~m of the saddle in at least \rev{99.6}\% of trials from every start up to 0.7~m (Fig.~\ref{fig:nonlinear}(b)). %DIF <  [REVIEW] hyperparameter study v3 E4, E5 (report_numbers.txt)
}%DIFDELCMD < 

%DIFDELCMD < %%%
\DIFdel{The attraction law equals $-k\hat{\mathbf{J}}^{-1}\hat{\mathbf{v}}(\mathbf{p}_c)$, a Newton step toward the zero of the local fit, and in continuous time such steps cannot carry the cluster across a curve where $\det\mathbf{J} = 0$. The start therefore largely decides which critical point is reached (Fig. ~\ref{fig:nonlinear}(a)). In \rev{99.89}\% of gyre and rotational-saddle trials, the sign of $\det\hat{\mathbf{J}}$ at the start matched the critical point reached, negative for a saddle and positive for a center. In the gyrethe center ended inside the triangle in \rev{92.3}\% of trials from starts inside the diamond, and most starts outside it ended at the nearest corner saddle. Every start within 0.3~m ended within 0.05~m of the center, while the true critical point drew every start to the center. Most gyre orbits at $r_d \ge 0.4$~m reach the diamond and are lost. %DIF <  [REVIEW] 451,508 trials; nearest critical point 91.3%; outside the diamond 10.3% reach the center, 60.6% overall; next two 0.1 m bands 74% and 17%; orbit survival 44% and 4% at r_d 0.4 and 0.5 m; curved saddle triangle rule 97.8% at chi <= 0.5 and 72.1% rotational at chi = 1. Sources: hyperparameter study v3 E5 basins, start_distance_bins.txt, E1hw gyre orbits; revision_metrics/results/triangle_*.csv, gyre_diamond_split (5000 starts); side ablation 0.20 and 0.45 m moves attraction rates by at most 1 point; the orbit rule scales with the triangle (chi = 0 saddle fails below 0.3 m at side 0.20, passes all at 0.45)
}%DIFDELCMD < 

%DIFDELCMD < \begin{figure}[!t]
%DIFDELCMD <     \centering
%DIFDELCMD <     %%%
\DIFdelFL{\includegraphics[width=\columnwidth]{fig_nonlinear_flows.pdf}}%DIFDELCMD < \\[-1pt]
%DIFDELCMD <     %%%
\makebox[0.576\columnwidth]{%DIFDELCMD < \footnotesize %%%
\DIFdelFL{(a)}}%DIFAUXCMD
\makebox[0.324\columnwidth]{%DIFDELCMD < \footnotesize %%%
\DIFdelFL{(b)}}
    %DIFAUXCMD
%DIFDELCMD < \caption{%
{%DIFAUXCMD
\DIFdelFL{Nonlinear fields in simulation. (a) Gyre cell and (b) irrotational curved saddle with $\chi = 1$, with streamlines in gray and \rev{six} plane-fit attraction trajectories from hollow starts. Blue trajectories reach the origin, and orange ones a neighboring saddle (squares). The dashed diamond is $\det\mathbf{J} = 0$.}} %DIFAUXCMD
%DIF <  [REVIEW] hyperparameter study v3: E1hw saved trajectories
    %DIFDELCMD < \label{fig:nonlinear}
%DIFDELCMD < \end{figure}
%DIFDELCMD < 

%DIFDELCMD < %%%
\DIFdelend \section{Hardware Validation}


\subsection{Hardware Testbed and Field Representation}

The testbed used for validation is SCU's Robotic Systems Laboratory Decabot testbed \cite{52}, which has been used to validate multirobot adaptive navigation algorithms for both cluster control and swarm particle optimization techniques \cite{8,53}. Fig.~\ref{fig:fullbed} shows the complete testbed.

\begin{figure}[!t]
    \centering
    \includegraphics[width=0.62\columnwidth]{testbed_with_four.png}
    \caption{Experimental testbed: 1.6 m × 1.6 m printed vector field map, Decabot rovers, and OptiTrack system. Four Decabots are shown; the experiments here used three of this fleet.}
    \label{fig:fullbed}
    \end{figure}

The experimental platform used three omnidirectional Decabot rovers operating within a 1.6~m $\times$ 1.6~m workspace. Each Decabot features three omnidirectional wheels providing holonomic motion capability and a downward-facing TCS34725 RGB color sensor enclosed in a light-shrouded housing to minimize external lighting effects. The robots' RGB sensors are calibrated against a printed palette spanning hue $0$ to $2\pi$ and saturation $0$ to $1$ at 24 intervals in each dimension. Each robot records raw RGB and clear-channel counts over every palette cell. The readings are normalized and stored as a four-dimensional feature vector $(R,G,B,K)$. Two feedforward networks, each with two hyperbolic-tangent hidden layers, are trained to recover the target hue and saturation from this vector. The saturation network has a single output, while the hue network outputs the sine and cosine of hue to respect its circularity. Training pools data across robots so that a single model transfers to an uncalibrated unit without per-robot tuning. Calibration achieved hue RMSE of 0.17 radians (9.7$^\circ$, R$^2$=0.991) and saturation RMSE of 4.4\% (R$^2$=0.951) across the operational range. The encoding and calibration palette follows~\cite{22}. All robots were calibrated after receiving preventive maintenance and servicing for this study.

The adaptive navigation and cluster space controller layers of Section~III run in MATLAB/Simulink on a central computer that communicates with the robots over Wi-Fi at 10~Hz. At each cycle it receives OptiTrack positions and RGB readings, computes the control laws, and returns velocity commands, which each robot's onboard computer (Intel Edison) converts into motor commands. The hardware controller also limited the commanded centroid velocity to 0.3~m/s per axis, a safeguard that reserves actuator capacity for formation keeping.

Vector fields are physically realized as printed floor maps using HSV color encoding, where hue (0--2$\pi$) represents flow direction and saturation (0--1) encodes magnitude. Printed magnitudes are scaled to the 0.3--1.0 saturation range to avoid undersaturated areas and ensure reliable sensor readings, then remapped to 0--1 in software, with saturation 1.0 equal to a field speed of 0.3~m/s. Together with the sensor shroud and self-illumination, this encoding is minimally affected by overhead lighting.

This method of representing vector fields in printed floor maps has imperfections in both printing and sensing. Fig.~\ref{fig:measurement_error} shows the analytical vortex and saddle fields, reconstructions of each field from sensed robot measurements, and the reconstruction error (right column). These reconstructed fields were used to fit multilayer perceptron (MLP) interpolation models, used in Section~IV as a fixed distortion model for simulation.


\begin{figure}[!t]
    \centering
    \includegraphics[width=\columnwidth]{fig_reconstruction_3x2.png}
    \caption{Analytical fields, sensor-based reconstructions, and reconstruction error for vortex (top) and saddle (bottom) fields.} % submitted measurement_error_comparison.png redrawn 3x2 from the same robot readings
    \label{fig:measurement_error}
\end{figure}


\subsection{Validation Test Plan}



Hardware validation focused on two representative field types that encompass the key challenges of vector field navigation: rotational dynamics (vortex) and hyperbolic instability (saddle). The vortex field tests pure rotational dynamics and provides continuity with previous studies in our group~\cite{22,48}. The saddle is the only canonical field type with separatrices, which act as Lagrangian coherent structures in steady flows. Testing on 
the saddle field validates that the framework operates reliably in the presence of these structures.

For each field, eight starting locations were selected around the workspace perimeter and the attraction controller (\ref{eq:10}) was run 10 times from each location for \open{20}~s using the three-robot equilateral triangle formation with a \open{0.33 or 0.30}~m side, yielding 80 trials per field. Run order was randomized. Every trial is scored with the metrics of Section~IV, and its final centroid speed is reported so that a cluster still moving is not counted as settled. The classifier of Section~IV-C labels $\hat{\mathbf{J}}$ on every cycle without feeding the controller, and its accuracy is scored against the type of the printed field. % [OPEN] run length, formation side (CAMPAIGN_PROMPT 4.2, 6.2)

Additionally, each field was used to test the critical point orbiting controller (\ref{eq:11}). A total of \open{six} commanded radii were chosen to match those in simulation, and each was run \open{three} times. Each orbiting controller experiment was run for \open{1800} cycles at 10~Hz, allowing the cluster to complete several orbits around the critical point. % [OPEN] number of radii, repeats, orbit run length (CAMPAIGN_PROMPT 6.2)

 \subsection{Attraction Controller Results}

Following the test plan above, \syn{80} vortex trials and \syn{78} saddle trials were analyzed. The final centroid position of each trial was used to determine systematic bias and precision, matching the simulation metrics in Section~IV. % [SYNTHETIC] run counts of the old 3R logs, table_IV_revised.tex; replaced by the campaign

The cluster converged to the critical point in \syn{100}\% of the saddle trials and \syn{86}\% of the vortex trials at the 0.05~m threshold, and in \syn{77}\% and \syn{34}\% at 0.03~m. No trial left the map, and the critical point ended inside the robots' triangle in \syn{100}\% of trials. Table~\ref{tab:estimation_results} summarizes the results. In the vortex, systematic bias was \syn{0.018}~m with precision of \syn{0.036}~m. In the saddle field, systematic bias was \syn{0.007}~m with precision of \syn{0.024}~m. The bias is smaller than the mean final error because final positions on opposite sides of $\mathbf{p}^*$ cancel in the mean, and the estimate ended closer to the critical point than the cluster. The type of $\hat{\mathbf{J}}$ matched the printed field on \open{campaign}\% of cycles in the vortex and \open{campaign}\% in the saddle. % [SYNTHETIC] old 3R logs scored by revision_metrics (table_IV_revised.tex, V_VIA numbers.tex), replaced by the campaign; [OPEN] hardware classification

Fig.~\ref{fig:vortextomiddle} shows convergence trajectories from all eight starting positions and mean centroid distance to the critical point over time for the same trials.

% [PLACEHOLDER] submitted Figs. 8 and 9, reused until the campaign. (a, b): allrunsadjusted.png split into
% its two panels; the vortex panel shows four-robot runs. (c): fig9.png with its title cropped; its vortex
% curve is the four-robot set (n=79) and its distances are in map units (OptiTrack / 1.6764).
% Replace all three panels with campaign data in metres.
\begin{figure}[!t]
    \centering
    \begin{minipage}[b]{0.47\columnwidth}\centering\includegraphics[width=\textwidth]{placeholder_hw_trajectories_a.png}\end{minipage}\hspace{0.04\columnwidth}%
    \begin{minipage}[b]{0.47\columnwidth}\centering\includegraphics[width=\textwidth]{placeholder_hw_trajectories_b.png}\end{minipage}\\[2pt]
    \includegraphics[width=0.92\columnwidth]{placeholder_hw_distance.png}\\[-1pt]
    {\footnotesize (c)}
    \caption{\syn{[Placeholder]} Convergence trajectories from eight starting positions to the center of the vortex (a) and saddle (b), with the solid line showing the mean path and the shaded region one standard deviation across the 10 repetitions. (c) Distance to center over time, with mean and standard deviation bounds.}
    \label{fig:vortextomiddle}\label{fig:Tomiddlewithtime}
\end{figure}

\begin{table}[!t]
    \footnotesize
    \centering
    \caption{Hardware experimental results \syn{(placeholder data)}. (a) Attraction, vortex: \syn{80} runs, saddle: \syn{78} runs. Bias measures mean offset from origin; precision ($\sigma$) measures repeatability. (b) Orbital control, radial errors in metres after the first 10~s. $\bar e_r$ is the radial bias, $s_b$ the between-trial spread, and $s_w$ the within-run standard deviation. Negative error indicates the cluster orbited inside the commanded radius.} % [SYNTHETIC] table_IV_revised.tex run counts (old 3R logs); V_VIA results/tab_orbit_body.tex
    \label{tab:estimation_results}\label{tab:orbital_radius_comparison}
    \setlength{\tabcolsep}{4pt}
    \syn{\begin{tabular}{l|ccc|ccc}
    \hline
     & \multicolumn{3}{c|}{\textbf{Vortex (printed map)}} & \multicolumn{3}{c}{\textbf{Saddle (printed map)}} \\
    \hline
    \multicolumn{7}{l}{\textit{(a) Attraction}} \\
    \hline
    Bias (m) & \multicolumn{3}{c|}{0.018} & \multicolumn{3}{c}{0.007} \\
    $\sigma$ (m) & \multicolumn{3}{c|}{0.036} & \multicolumn{3}{c}{0.024} \\
    $\bar{e}_f \pm s_e$ (m) & \multicolumn{3}{c|}{0.037 $\pm$ 0.017} & \multicolumn{3}{c}{0.023 $\pm$ 0.009} \\
    Median [p95] (m) & \multicolumn{3}{c|}{0.036 [0.070]} & \multicolumn{3}{c}{0.024 [0.037]} \\
    Loc.\ error (m) & \multicolumn{3}{c|}{0.028} & \multicolumn{3}{c}{0.018} \\
    $T_{\mathrm{conv}}$ (s) & \multicolumn{3}{c|}{3.6} & \multicolumn{3}{c}{3.3} \\
    $E_u$ (m$^2$/s) & \multicolumn{3}{c|}{0.187} & \multicolumn{3}{c}{0.205} \\
    Succ.\ 3/4/5 cm (\%) & \multicolumn{3}{c|}{34/60/86} & \multicolumn{3}{c}{77/99/100} \\
    \hline
    \multicolumn{7}{l}{\textit{(b) Orbit}} \\
    \hline
    $r_d$ (m) & $\bar e_r$ & $s_b$ & $s_w$ & $\bar e_r$ & $s_b$ & $s_w$ \\
    \hline
    0.01 & +0.166 & 0.031 & 0.027 & +0.111 & 0.014 & 0.018 \\
    0.10 & +0.151 & 0.028 & 0.061 & +0.074 & 0.015 & 0.028 \\
    0.20 & +0.128 & 0.027 & 0.064 & +0.071 & 0.018 & 0.047 \\
    0.30 & +0.112 & 0.027 & 0.097 & +0.091 & 0.033 & 0.059 \\
    0.40 & +0.081 & 0.036 & 0.119 & +0.102 & 0.040 & 0.079 \\
    0.50 & +0.018 & 0.020 & 0.079 & +0.088 & 0.041 & 0.095 \\
    \hline
    \end{tabular}} % [SYNTHETIC] revision_metrics/tables/table_IV_revised.tex (old 3R logs) and V_VIA results/tab_orbit_body.tex (three-robot columns, rows 0.05 and 0.15 m omitted), replaced by the campaign
    \end{table}



\subsection{Orbital Control Results}

Orbital control was evaluated on both fields at commanded radii from 0.01~m to 0.50~m. For each radius, the orbital controller (\ref{eq:11}) was run for \open{1800} cycles at 10~Hz, and the mean radial error and its standard deviation were recorded after the first 10~s, as in simulation (Table~\ref{tab:orbital_radius_comparison}). % [OPEN] orbit run length (CAMPAIGN_PROMPT 6.2)

 In the vortex field, mean radial errors decreased from \syn{0.166}~m at a commanded radius of 0.01~m to \syn{0.018}~m at 0.50~m, with within-run standard deviations ranging from \syn{0.027} to \syn{0.119}~m. In the saddle field, mean radial errors ranged from \syn{0.111}~m to \syn{0.088}~m over the same commanded radii, with within-run standard deviations increasing from \syn{0.018} to \syn{0.095}~m. From $r_d = \syn{0.10}$~m upward the fluctuation within one orbit exceeds the variation between repeats. By the $0.25\,r_d$ criterion the practical range is $r_d \ge \syn{0.40}$~m in the vortex and $r_d \ge \syn{0.50}$~m in the saddle. Below these radii the controller holds a bounded orbit about the critical point but does not track the commanded radius. \syn{One} vortex and \syn{three} saddle orbits failed by leaving the map, all at $r_d \ge \syn{0.40}$~m. % [SYNTHETIC] V_VIA numbers.tex and tab_orbit_body.tex (three-robot columns, six radii)




\section{Discussion}

\subsection{Error Analysis}

Estimation and control errors arise from four sources: robot dynamics, field distortion and measurement noise, cluster geometry, and field conditioning. Sections~II-B and~III-D bound the effect of each, and this section traces the errors observed in Sections~IV and~V to them.

Perfect convergence does not occur in simulation. However, repeating experiments with no stiction, no max velocity, and $\alpha = 0$ (instantaneous velocity response) confirms convergence to the controller's $10^{-6}$~m command deadband, with the estimate $\hat{\mathbf{p}}^*$ itself matching $\mathbf{p}^*$ to machine precision. This confirms the variance is due to robot dynamics, not the estimation method. Some errors may also have arisen from modeling discrepancies such as modeling all robots identically with single stiction, max speed, and response time values.

Positive radial bias at small orbital radii is also a result of robot dynamics. At small radii, the commanded angular rate $k_t / r_d$ is high relative to the radial correction bandwidth $k_r$, and the lag-induced drift of Section~III-D sets the bias. The radial term of (\ref{eq:11}) cancels that drift when $k_r(r^* - r_d) = \tau V^2/r^*$, whose positive root is
\begin{equation}\label{eq:orbit_cf}
r^* = \tfrac{1}{2}\bigl(r_d + \sqrt{r_d^2 + 4C}\bigr), \qquad C = \tau V^2/k_r.
\end{equation}
The error $r^* - r_d$ approaches $r_{\min} = \sqrt{C}$ as $r_d \to 0$ and falls as $C/r_d$ at large radii. Using the effective lag of the sampled loop in place of $\tau$, the closed form matches simulation to a median of \rev{0.03}~mm across sweeps of gains, lag, sample period, and friction. At the settings of Section~IV, $r_{\min} = \rev{0.133}$~m, and the error falls below 0.05~m only for $r_d \ge \rev{0.31}$~m. Lowering $k_t$ to 0.1~m/s cuts $r_{\min}$ to \rev{0.058}~m, raising $k_r$ to 2~s$^{-1}$ cuts it to \rev{0.067}~m, and removing the lag cuts it to \rev{0.059}~m. % [REVIEW] hyperparameter study v3 report_numbers.txt (closed form, table 3)
A higher radial gain also raises the sensitivity to estimate noise.

Comparing the noise-free and reconstructed fields in simulation isolates the effect of field distortion: the fixed distortion of the printed maps appears to be the driving factor of the simulation-reality gap, with its errors propagating directly through the critical point estimation equations. A parked robot's reading changes by about \syn{$2\times10^{-5}$}~m/s between cycles, while the calibration error of Section V-A is about 0.02~m/s and is fixed at each location. % [SYNTHETIC] old hardware logs (prose batch N1); campaign block E replaces it
In the final error of Section~III-D, friction sets the spread of final positions through $v_c/k$, and the distortion sets their bias through the estimate error at the critical point.

The hardware orbit metrics separate estimation error from control error. At $r_d = 0.01$~m the controller-only bias in the vortex was \syn{0.177}~m against a true bias of \syn{0.166}~m, so the small-radius offset arises in the closed loop. The localization error grows with distance from the critical point, and the within-run spread grows with it. Runs with the formation parked at fixed points separate temporally varying noise from the fixed distortion \open{campaign block E}. % [SYNTHETIC] V_VIA numbers.tex; [OPEN] static noise runs

To avoid errors from collinear robots (the condition number of $\mathbf{A}$ becomes infinite), an equilateral triangle was chosen as the cluster shape. It minimizes the condition number of $\mathbf{A}$ in (\ref{eq:4}) among formations of the same size centered on the sampled region, making it optimal for sampling.

Further errors could arise from poor conditioning of the estimated Jacobian $\kappa(\hat{\mathbf{J}})$ in (\ref{eq:8}), directly affecting the critical point estimate (\ref{eq:9}). The drop of the orbital bias below the controller-only bias at the largest commanded radius (Table~\ref{tab:simulation_orbital_results}) traces to degradation of the reconstructed field near the edge of the mapped region: there the estimated critical point distance is systematically inflated, so the radial controller pulls the cluster inward. This is the range-dependent error amplification captured by $\kappa(\hat{\mathbf{J}})$, which measures how the local field structure amplifies coefficient errors into location errors. Unlike the condition number of $\mathbf{A}$, this factor is an intrinsic property of the field and is not under the designer's control. In the limiting case, $\kappa(\hat{\mathbf{J}}) \to \infty$ as $\det(\hat{\mathbf{J}}) \to 0$, corresponding to degenerate critical points or near-uniform flow where no isolated critical point exists; the cluster approaches this condition far from the critical point, where the sampled field looks locally uniform. These singular cases are detectable at run time through eigenvalue analysis of the estimated Jacobian, and Section~III-D describes how the controller handles them. A center is non-hyperbolic yet has $\det\mathbf{J} > 0$, so (\ref{eq:9}) locates it as well as any other critical point, and only its type is fragile (Section~IV-C).
\subsection{Comparison with Matched Baselines}

To separate the estimate from the controller, the attraction law (\ref{eq:10}) and the orbital law (\ref{eq:11}) were run in simulation with the same gains, robot model, and formation, driven by three critical point estimates. The oracle uses the true critical point and isolates the controller. The magnitude estimate uses only the sensed speeds $S_i = \|\mathbf{m}_i\|$, as in \cite{48}. A plane fit gives their slope $\mathbf{g}$ and their value $S_c$ at the centroid, and $\hat{\mathbf{p}}_m = \mathbf{p}_c - S_c\mathbf{g}/\|\mathbf{g}\|^2$ locates the zero of a speed that grows linearly with distance. It uses the gains of the proposed laws without separate tuning. Table~\ref{tab:baselines} reports both beside the plane fit.

On the affine fields the plane fit matches the oracle in every metric (Table~\ref{tab:baselines}, Fig.~\ref{fig:effort}), because its estimate is exact, and the remaining errors belong to the controller. % [REVIEW] vib_baselines_note, revision_metrics/tables/table_baselines_revised.tex

The magnitude estimate errs most near the critical point, where the sensed speeds form a cone and the fitted slope vanishes. Its cluster never comes to rest and still moves at about \rev{0.12}~m/s after 10~s. Its smaller final error is therefore a snapshot of that motion, bought with about \rev{four} times the control effort (Fig.~\ref{fig:effort}). In orbit, an inward estimate bias offsets the lag at 0.3 and 0.5~m on the affine fields but holds the orbit \rev{0.07}~m inside the commanded 0.1~m. The canonical sink, source, vortex, and saddle share the speed field $\|\mathbf{v}\| = r$, so only the plane fit identifies the type of the critical point. % [REVIEW] vib_baselines_note

\begin{figure}[!t]
    \centering
    \includegraphics[width=\columnwidth]{fig_sim_speed.pdf}
    \caption{Centroid speed in simulated attraction with three estimates, from \rev{three} starts of Table~\ref{tab:baselines}, \rev{0.49 to 0.56}~m from the critical point. Control effort $E_u$ is the area under the squared speed. The plane fit matches the oracle on the analytical vortex, and the magnitude estimate never brings the cluster to rest.} % [REVIEW] run_speed_traces.py: Table VI attraction batches rerun with recording (seed 0, Coulomb model); effort and final error match baselines_attraction.csv; traces in revision_metrics/results/baselines_speed_traces.npz
    \label{fig:effort}
\end{figure}

\begin{table}[!t]
\centering
\scriptsize
\caption{Matched baselines in simulation, \rev{500} attraction trials of 10~s and \rev{50} orbit trials of 60~s per field and setting. Success is ending within 0.05~m, orbit RMSE of $r^* - r_d$ excludes the first 10~s, and Est.\ err.\ is the mean $\|\hat{\mathbf{p}}^* - \mathbf{p}^*\|$ during the orbits. Vortex and saddle are reconstructed, and the oracle gives the same values on every field.} % [REVIEW] revision_metrics/tables/table_baselines_revised.tex caption (trial counts)
\label{tab:baselines}
\setlength{\tabcolsep}{2pt}
\rev{\begin{tabular}{lcccc|ccc|c}
\hline
 & \multicolumn{4}{c|}{\textbf{Attraction}} & \multicolumn{3}{c|}{\textbf{Orbit RMSE (m)}} & \textbf{Est.} \\
\textbf{Field,} & \textbf{Succ.} & \textbf{$\bar{e}_f$} & \textbf{$T_{\mathrm{conv}}$} & \textbf{$E_u$} & \multicolumn{3}{c|}{\textbf{at $r_d$ (m)}} & \textbf{err.} \\
\textbf{estimate} & \textbf{(\%)} & \textbf{(m)} & \textbf{(s)} & \textbf{(m$^2$/s)} & 0.1 & 0.3 & 0.5 & \textbf{(m)} \\
\hline
All, oracle & 100 & 0.025 & 2.1 & 0.074 & 0.093 & 0.051 & 0.033 & 0.000 \\
Analyt., magn. & 100 & 0.012 & 1.3 & 0.306 & 0.068 & 0.016 & 0.020 & 0.060 \\
Vortex, fit & 94 & 0.031 & 2.3 & 0.069 & 0.096 & 0.092 & 0.045 & 0.051 \\
Vortex, magn. & 94 & 0.034 & 1.8 & 0.269 & 0.065 & 0.079 & 0.068 & 0.211 \\
Saddle, fit & 100 & 0.025 & 2.3 & 0.071 & 0.090 & 0.071 & 0.050 & 0.061 \\
Saddle, magn. & 100 & 0.018 & 1.5 & 0.270 & 0.066 & 0.061 & 0.056 & 0.214 \\
\hline
\end{tabular}} % [REVIEW] revision_metrics/tables/table_baselines_revised.tex; identical oracle rows and the analytical plane-fit row merged
\end{table}


\subsection{Limitations}

The estimator fits a planar affine model across the formation, which limits the fields it can address. It applies only to planar vector fields, and a 3D field would require at least four non-coplanar robots for Jacobian estimation. It also assumes the affine model holds from the cluster out to the critical point (Section~II-B). On fields that are highly nonlinear at the scale of the formation, the fit degrades and biases the recovered critical point location. Section~IV\DIFdelbegin \DIFdel{-D}\DIFdelend  probes this regime on a gyre cell and a curved saddle, but fields with stronger curvature, or with critical points closer together than one formation side, were not tested.

All experiments used three robots, the minimum. With $n > 3$ robots the least-squares fit averages measurement errors across robots and survives the loss of one. However, a wider footprint asks the affine model to hold over more of the field's curvature, and each added robot brings more communication load and two more shape variables for the formation controller to regulate.

The printed fields give exact ground truth, but they are static and do not push the robots. A physical flow advects the cluster and adds its own velocity $\mathbf{J}\mathbf{e}$ to the centroid near $\mathbf{p}^*$. For an ideal centroid with an exact estimate, the error then obeys $\dot{\mathbf{e}} = (\mathbf{J} - k\mathbf{I})\mathbf{e}$, which is unstable when an eigenvalue of $\mathbf{J}$ has real part above $k$. Actuator lag lowers this limit in rotating flow. A physical critical point can also drift. Even without advection, the attraction law trails one moving at constant speed by that speed divided by $k$.

The computation is centralized. A central computer receives OptiTrack positions and field readings from every robot, forms the estimate, and returns velocity commands. The estimate needs only the positions and readings of one cycle in a common frame, so a robot receiving them from its neighbors could compute it. Onboard relative localization and inter-robot communication were not tested. In either architecture, every fit depends on the communication that delivers its readings. A dropout can leave fewer than three readings, and the fit then has no unique solution. Latency mixes readings from different times into one fit.

\subsection{Future Work}

The estimated Jacobian's eigenvalues determine the topological type of the nearby critical point (Table~\ref{tab:field_summary}), and Sections~IV-C and V-C evaluate this classification with a fixed band around the center and node boundaries. Replacing the band with a test against the covariance of $\hat{\mathbf{J}}$ under measurement noise would label fragile types with a stated confidence. \DIFaddbegin \DIFadd{Because the sign of $\det\hat{\mathbf{J}}$ on the first cycle decided the type of critical point reached (Section~IV-A), the same test could hold the cluster in place when the sensed type differs from the one the mission seeks.
}\DIFaddend 

Because the three-robot affine fit developed here returns a spatially constant Jacobian $\hat{\mathbf{J}}$, it cannot resolve how the velocity gradient varies across the formation. Extending the estimator to a second-order polynomial fit from six robots that do not lie on a common conic recovers a position-dependent Jacobian across the footprint. Work is being done to see how these new estimates can form controllers that can detect other topological features of practical interest beyond critical points.



\section{Conclusion}
\noindent

This paper presents a framework for multirobot navigation in vector fields containing critical points, in which the robots sense the field at spatially distributed points and a central computer forms the estimate. A Jacobian-based approach uses instantaneous measurements from three robots to identify critical point locations without prior field knowledge, supporting attraction and orbital control laws built on these estimates. Simulation across six canonical field types confirmed convergence of these control primitives. On curved fields the estimate degraded as the cube of the formation side, and the \DIFdelbegin \DIFdel{cluster reached the critical point whose basin, bounded by $\det\mathbf{J} = 0$, contained its start}\DIFdelend \DIFaddbegin \DIFadd{sign of $\det\hat{\mathbf{J}}$ on the first cycle decided which type of critical point the cluster reached}\DIFaddend . In \syn{158} hardware convergence trials, the system achieved a \syn{100}\% (saddle) and \syn{86}\% (vortex) success rate within 0.05~m, with mean position errors of \syn{0.037}~m or less. Orbital control experiments maintained commanded radii within a quarter of the radius from \syn{0.40}~m (vortex) and \syn{0.50}~m (saddle) upward, and below that held bounded orbits that did not track the commanded radius. % [SYNTHETIC] hardware numbers
The printed fields do not advect the robots, so these results hold under lab conditions. A physical flow adds a disturbance to the robot dynamics that the present controller does not compensate for.


\section*{Disclosure of Conflicts of Interest}

C. A. Kitts serves as an Associate Editor of the \textit{IEEE Systems Journal}. He has been excluded from all editorial handling and peer review of this article, including the assignment of an Associate Editor and the selection of reviewers, and has been designated a non-preferred Associate Editor in the submission system.

C. Waight and C. A. Kitts are co-inventors on a provisional patent application filed through Santa Clara University covering the method described in this article. Santa Clara University holds any rights arising from that application.

The authors declare no other financial or non-financial competing interests.


%\balance % off while the last page has room; restore once the enlarged figures fill it
\begin{thebibliography}{99}

\bibitem{1} C. Kitts and M. Egerstedt, "Design, control, and applications of real-world multirobot systems [from the guest editors]," \textit{IEEE Robot. Autom. Mag.}, vol. 15, no. 1, p. 8, Mar. 2008.

\bibitem{2} J. Cortés and M. Egerstedt, "Coordinated control of multi-robot systems: A survey," \textit{SICE J. Control Meas. Syst. Integr.}, vol. 10, no. 6, pp. 495--503, 2017.

\bibitem{3} P. Ögren, E. Fiorelli, and N. E. Leonard, "Cooperative control of mobile sensor networks: Adaptive gradient climbing in a distributed environment," \textit{IEEE Trans. Autom. Control}, vol. 49, no. 8, pp. 1292--1302, 2004.

\bibitem{4} L. Briñón-Arranz, A. Renzaglia, and L. Schenato, "Multirobot symmetric formations for gradient and hessian estimation with application to source seeking," \textit{IEEE Trans. Robot.}, vol. 35, no. 3, pp. 782--789, 2019.

\bibitem{5} C. A. Kitts, R. T. McDonald, and M. A. Neumann, "Adaptive navigation control primitives for multirobot clusters: Extrema finding, contour following, ridge/trench following, and saddle point station keeping," \textit{IEEE Access}, vol. 6, pp. 17625--17642, 2018.

\bibitem{6} R. T. McDonald, C. A. Kitts, and M. A. Neumann, "Experimental implementation and verification of scalar field ridge, trench, and saddle point maneuvers using multirobot adaptive navigation," \textit{IEEE Access}, vol. 7, pp. 62950--62961, 2019.

\bibitem{7} R. K. Lee et al., "Multiple UAV adaptive navigation for three-dimensional scalar fields," \textit{IEEE Access}, vol. 9, pp. 122626--122654, 2021.

\bibitem{8} R. T. McDonald, C. A. Kitts, and M. A. Neumann, "Navigation of scalar fronts with multirobot clusters in simulation and experiment," \textit{IEEE Syst. J.}, vol. 14, no. 3, pp. 3755--3766, 2020.

\bibitem{9} R. K. Lee, C. A. Kitts, M. A. Neumann, and R. T. McDonald, "3-d adaptive navigation: Multirobot formation control for seeking and tracking of a moving source," in \textit{Proc. IEEE Aerosp. Conf.}, pp. 1--12, 2020.

\bibitem{10} T. Adamek, C. A. Kitts, and I. Mas, "Gradient-based cluster space navigation for autonomous surface vessels," \textit{IEEE/ASME Trans. Mechatronics}, vol. 20, no. 2, pp. 506--518, 2014.

\bibitem{11} F. Pagano, N. De Carli, E. Restrepo, A. Marino, and P. Robuffo Giordano, "Distributed multi-robot active-sensing of a diffusive source," \textit{IEEE Robot. Autom. Lett.}, vol. 10, no. 10, pp. 10807--10814, 2025.

\bibitem{13} S. Li, R. Kong, and Y. Guo, "Cooperative distributed source seeking by multiple robots: Algorithms and experiments," \textit{IEEE/ASME Trans. Mechatronics}, vol. 19, no. 6, pp. 1810--1820, 2014.

\bibitem{14} Z. Deng and J. Luo, "Fully distributed algorithms for constrained nonsmooth optimization problems of general linear multiagent systems and their application," \textit{IEEE Trans. Autom. Control}, vol. 69, no. 2, pp. 1377--1384, 2023.

\bibitem{15} W. Li, J. A. Farrell, S. Pang, and R. M. Arrieta, "Moth-inspired chemical plume tracing on an autonomous underwater vehicle," \textit{IEEE Trans. Robot.}, vol. 22, no. 2, pp. 292--307, 2006.

\bibitem{18} S. R. Ramp et al., "Preparing to predict: the second autonomous ocean sampling network (AOSN-II) experiment in the Monterey Bay," \textit{Deep Sea Res. Part II: Top. Stud. Oceanogr.}, vol. 56, no. 3-5, pp. 68--86, 2009.

\bibitem{19} J. Das et al., "Coordinated sampling of dynamic oceanographic features with underwater vehicles and drifters," \textit{Int. J. Robot. Res.}, vol. 31, no. 5, pp. 626--646, 2012.

\bibitem{20} J. Wang et al., "Dynamic plume tracking by cooperative robots," \textit{IEEE/ASME Trans. Mechatronics}, vol. 24, no. 2, pp. 609--620, 2019.

\bibitem{21} K. Garg and D. Panagou, "Finite-time estimation and control for multi-aircraft systems under wind and dynamic obstacles," \textit{J. Guid. Control Dyn.}, vol. 42, no. 7, pp. 1489--1505, 2019.

\bibitem{22} C. Waight and C. Kitts, "A functional indoor testbed for multirobot adaptive navigation in vector field environments," in \textit{Proc. ASME Int. Design Eng. Tech. Conf. Comput. Inf. Eng. Conf.}, vol. 89251, 2025.

\bibitem{23} M. Michini, M. A. Hsieh, E. Forgoston, and I. B. Schwartz, "Robotic tracking of coherent structures in flows," \textit{IEEE Trans. Robot.}, vol. 30, no. 3, pp. 593--603, 2014.

\bibitem{24} D. R. Nelson, D. B. Barber, T. W. McLain, and R. W. Beard, "Vector field path following for miniature air vehicles," \textit{IEEE Trans. Robot.}, vol. 23, no. 3, pp. 519--529, 2007.

\bibitem{25} W. Yao, H. G. de Marina, B. Lin, and M. Cao, "Singularity-free guiding vector field for robot navigation," \textit{IEEE Trans. Robot.}, vol. 37, no. 4, pp. 1206--1221, 2021.

\bibitem{26} D. Panagou, "Motion planning and collision avoidance using navigation vector fields," in \textit{Proc. IEEE Int. Conf. Robot. Autom.}, 2014, pp. 2513--2518.

\bibitem{28} X. Wang et al., "Adaptive vector field guidance without a priori knowledge of course dynamics and wind," \textit{IEEE/ASME Trans. Mechatronics}, vol. 27, no. 6, pp. 4597--4607, 2022.

\bibitem{29} Y. Zhang, M. A. Godin, J. G. Bellingham, and J. P. Ryan, "Using an autonomous underwater vehicle to track a coastal upwelling front," \textit{IEEE J. Ocean. Eng.}, vol. 37, no. 3, pp. 338--347, 2012.

\bibitem{32} D. Kularatne and M. A. Hsieh, "Tracking attracting manifolds in flows," \textit{Auton. Robots}, vol. 41, no. 8, pp. 1575--1588, 2017.

\bibitem{36} X. Tricoche, G. Scheuermann, and H. Hagen, "Continuous topology simplification of planar vector fields," in \textit{Proc. IEEE Vis.}, pp. 159--166, 2001.

\bibitem{37} W. G. Pichel et al., "Marine debris collects within the North Pacific subtropical convergence zone," \textit{Mar. Pollut. Bull.}, vol. 54, no. 8, pp. 1207--1211, 2007.

\bibitem{38} M. Mateus, R. Canelas, L. Pinto, and N. Vaz, "When tragedy strikes: potential contributions from ocean observation to search and rescue operations after drowning accidents," \textit{Front. Mar. Sci.}, vol. 7, p. 55, 2020.

\bibitem{40} P. Keramea, K. Spanoudaki, G. Zodiatis, G. Gikas, and G. Sylaios, "Oil spill modeling: A critical review on current trends, perspectives, and challenges," \textit{J. Mar. Sci. Eng.}, vol. 9, no. 2, p. 181, 2021.

\bibitem{41} J. Wang, Y. Lin, R. Liu, and J. Fu, "Odor source localization of multi-robots with swarm intelligence algorithms: A review," \textit{Front. Neurorobot.}, vol. 16, p. 949888, 2022.

\bibitem{43} C. Dong et al., "Oceanic mesoscale eddies," \textit{Ocean-Land-Atmos. Res.}, vol. 4, p. 0081, 2025.

\bibitem{44} Z. Li, X. Wang, J. Hu, F. P. Andutta, and Z. Liu, "A study on an anticyclonic-cyclonic eddy pair off Fraser Island, Australia," \textit{Front. Mar. Sci.}, vol. 7, p. 594358, 2020.

\bibitem{45} R. Deogharia, H. Gupta, S. Sil, A. Gangopadhyay, and A. Shee, "On the evidence of helico-spiralling recirculation within coherent cores of eddies using Lagrangian approach," \textit{Sci. Rep.}, vol. 14, no. 1, p. 11014, 2024.

\bibitem{30} T. Inanc, S. C. Shadden, and J. E. Marsden, "Optimal trajectory generation in ocean flows," in \textit{Proc. Amer. Control Conf.}, pp. 674--679, 2005.

\bibitem{31} R. N. Smith, M. Schwager, S. L. Smith, B. H. Jones, D. Rus, and G. S. Sukhatme, "Persistent ocean monitoring with underwater gliders: Adapting sampling resolution," \textit{J. Field Robot.}, vol. 28, no. 5, pp. 714--741, 2011.

\bibitem{33} G. Knizhnik, P. Li, X. Yu, and M. A. Hsieh, "Flow-based control of marine robots in gyre-like environments," in \textit{Proc. IEEE Int. Conf. Robot. Autom. (ICRA)}, pp. 3047--3053, 2022.

\bibitem{46} Y. Jiao, H. Hang, J. Merel, and E. Kanso, "Sensing flow gradients is necessary for learning autonomous underwater navigation," \textit{Nat. Commun.}, vol. 16, no. 1, p. 3044, 2025.

\bibitem{47} D. Rajabhandharaks, R. T. McDonald, M. A. Neumann, and C. A. Kitts, "Indoor testbed for vector field multirobot adaptive navigation: Feasibility study," in \textit{Proc. ASME Int. Design Eng. Tech. Conf. Comput. Inf. Eng. Conf.}, vol. 59292, p. V009T12A043, 2019.

\bibitem{48} R. Cooper et al., "Initial Study of Multirobot Adaptive Navigation for Exploring Environmental Vector Fields," in \textit{Proc. IEEE Aerosp. Conf.}, 2019.

\bibitem{49} Z. Zhou, J. Liu, and J. Yu, "A survey of underwater multi-robot systems," \textit{IEEE/CAA J. Autom. Sinica}, vol. 9, no. 1, pp. 1--18, 2021.

\bibitem{50} C. A. Kitts and I. Mas, "Cluster space specification and control of mobile multirobot systems," \textit{IEEE/ASME Trans. Mechatronics}, vol. 14, no. 2, pp. 207--218, 2009.

\bibitem{51} I. Mas and C. Kitts, "Obstacle avoidance policies for cluster space control of nonholonomic multirobot systems," \textit{IEEE/ASME Trans. Mechatronics}, vol. 17, no. 6, pp. 1068--1079, 2012.

\bibitem{52} S. Tomer et al., "A low-cost indoor testbed for multirobot adaptive navigation research," in \textit{Proc. IEEE Aerosp. Conf.}, pp. 1--12, 2018.

\bibitem{53} S. T. Hart and C. A. Kitts, "Unifying control architecture for reactive particle swarms," \textit{IEEE/ASME Trans. Mechatronics}, vol. 28, no. 2, pp. 873--883, 2022.

\end{thebibliography}

\begin{IEEEbiography}[{\includegraphics[width=1in,height=1.25in,clip,keepaspectratio]{fig1.jpg}}]{Christopher Waight}
received his BSc degree in Physics from Santa Clara University. He later earned his Master's in Mechanical Engineering there, and is now a PhD Candidate in Mechanical Engineering. He currently works as a Research Assistant in SCU's Robotic Systems Laboratory. He also has experience working as a hardware engineer and as a Technical Program Manager in R\&D.
\end{IEEEbiography}

\begin{IEEEbiography}[{\includegraphics[width=1in,height=1.25in,clip,keepaspectratio]{ckitts.png}}]{Christopher A. Kitts}
    (Senior Member, IEEE) earned degrees from Princeton University (B.S.E.), University of Colorado (M.P.A.), and Stanford University (M.S. and Ph.D.). He served as a Satellite Constellation Mission Controller with the U.S. Air Force and as a Computer Scientist at NASA Ames Research Center. He is currently Professor of Mechanical Engineering, Director of the Robotic Systems Laboratory, and Associate Dean at Santa Clara University. Dr. Kitts is a Fellow of ASME.
    \end{IEEEbiography}

\end{document}
